\documentclass[aspectratio=169,11pt,xcolor=table]{beamer}
\usepackage{slidestyle}
\ifdefined\notesmode
  \setbeameroption{show only notes}
\else
  \setbeameroption{hide notes}
\fi
\title{Green Machines + Metallic materials}

\begin{document}

% ================================================================ opening
\begin{frame}
\relax
\vspace*{4mm}
{\usebeamerfont{title}\usebeamercolor[fg]{frametitle}Green Machines + Metallic materials}\par
\vspace{1mm}{\color{rule}\rule{\textwidth}{0.8pt}}\par\vspace{2mm}
Year 9 Science \textperiodcentered{} exam sprint \textperiodcentered{} 2 hours\par\vspace{5mm}
\begin{tabular}{@{}l@{\hspace{6mm}}l@{}}
0:00 & Lesson 1 \textperiodcentered{} Green Machines\\
0:50 & Lesson 2 \textperiodcentered{} Metallic materials\\
1:40 & Mini mock (both units)\\
\end{tabular}
\zh{%
屏幕：课题 Green Machines（植物单元）+ Metallic materials（金属单元），Year 9 Science，考前冲刺，2 小时。下面是时间表：0:00 第一课 Green Machines；0:50 第二课 Metallic materials；1:40 Mini mock（小模考，两个单元都考）。\par
背景：这是两本学校 booklet 的复习课。Green Machines 讲植物怎么繁殖、怎么制造养分、怎么运输水分；Metallic materials 讲金属的性质、密度、合金、生锈和化学反应。考试是英文，评分用 A / M / E：A (Achieved) 写出正确的事实，M (Merit) 把原因连起来，E (Excellence) 整条因果链完整、并用数据或对比。}
\note{%
Teaching line, not a question slide.\par
两小时考前冲刺，一对一，投屏远距离。学生基础好，节奏可以快。\par
时间表：Lesson 1 Green Machines 0:00–0:50；Lesson 2 Metallic materials 0:50–1:40；Mini mock 1:40–1:57；作业 1:57–2:00。\par
开课前发纸：Classwork (Part A–F) 和 Mini mock 先发 Classwork，Mini mock 到 1:40 再发；Homework 两份下课时发。全部下课收回。\par
评分语言统一用 A / M / E：A 写出事实，M 把原因连上，E 把整条链写完整并用数据或对比。\par
课件页脚的页码指她的两本 booklet：Green Machines p.X 和 Metals p.X。\par
Script:\par
说：“今天两个小时，前 50 分钟复习 Green Machines (植物) ，后 50 分钟复习 Metallic materials (金属) ，最后 17 分钟做一份 mini mock (小模考) 。”\par
说：“每一段都会先快速回忆单词，再讲重点，再找 Sam 的错误，最后做几道练习题。”\par
说：“回答问题的时候，科学词一定用英文说，因为考试要写英文。”}
\end{frame}

% ================================================================ LESSON 1
\begin{frame}
\relax
\lessondivider{LESSON 1 \textperiodcentered{} Green Machines}{Learning intention: explain how a plant makes food, moves water and makes new plants, using the exam words.}
\zh{%
屏幕：LESSON 1 · Green Machines。学习目标：用考试用的词，解释植物怎么制造食物、怎么运输水分、怎么产生新的植物。\par
背景：Green Machines 就是“绿色机器”，指植物：它们自己用阳光制造养分。}
\note{%
Teaching line, not a question slide.\par
时间 0:00–0:02。\par
告诉她这一小时的三段：Loop 1 reproduction，Loop 2 photosynthesis and respiration，Loop 3 transport。每段结尾 2–3 道 Practice，最后 10 分钟做 Classwork Part A–C。\par
她这本 booklet 基本都写了，所以今天不是重新教，而是把答案写到 M 和 E。\par
Script:\par
说：“第一课，Green Machines。我们要解释三件事：plant 怎么 make food，怎么 move water，怎么 make new plants。”\par
说：“你 booklet 基本都做完了，所以今天的目标不是重新学，而是把答案写到 Merit 和 Excellence 的水平。”\par
Say it:\par
photosynthesis = foh-toh-SIN-thuh-sis\par
respiration = res-pih-RAY-shun}
\end{frame}

\begin{frame}{Do Now}
\relax
\qline{1} Which part of a flower makes pollen?
\qline{2} The pollen nucleus fuses with the ovule nucleus. Is this pollination or fertilisation?
\qline{3} Name the three conditions a seed needs to germinate.
\qline{4} Which gas does a leaf take in to make glucose?
\qline{5} Which tubes carry water up the stem?
\fref{Green Machines p.2}
\zh{%
屏幕：Do Now（课前小测）。1. 花的哪一部分产生 pollen (花粉)？2. 花粉的 nucleus (细胞核) 和 ovule (胚珠) 的细胞核融合，这叫 pollination (授粉) 还是 fertilisation (受精)？3. 种子 germinate (萌发) 需要哪三个条件？4. 叶子吸收哪种气体来制造 glucose (葡萄糖)？5. 哪种管道把水沿着茎往上运？\par
背景：Do Now 是一开始的 5 道快问快答，用来看她哪里还不熟。答案都在备注里，Diagnostic 一栏告诉你答错说明什么。}
\note{%
Q1 anther\par
Q2 fertilisation\par
Q3 water, warmth, oxygen\par
Q4 carbon dioxide\par
Q5 xylem\par
Short solution:\par
Q1 anther 在 filament 顶端，两者合起来是 stamen\par
Q2 pollen 从 anther 到 stigma 是 pollination；两个 nucleus 融合才叫 fertilisation\par
Watch for:\par
Q2 答 pollination 说明两个词混了，这是她最容易丢的定义分\par
Q5 答 “veins” 不够，要写 xylem\par
Diagnostic:\par
Q2 错：Loop 1 的 “From pollen to seed” 多花 1 分钟\par
Q4 答 oxygen：把 photosynthesis 和 respiration 的方程式记反了，Loop 2 先讲方程式\par
Q5 答 phloem：xylem 和 phloem 混淆，Loop 3 的 Quick review 先做\par
Script:\par
说：“先做 Do Now，五道题，每题一个词或一句话，用英文说答案。”\par
念第 1 题：“Which part of a flower makes pollen? 花的哪个部分产生 pollen？” 她应答 anther (花药)。\par
念第 2 题：“Pollination or fertilisation?” 她应答 fertilisation。如果她答 pollination，说：“pollination 只是把 pollen 送到 stigma (柱头) 上；两个 nucleus 合在一起才叫 fertilisation。”\par
念第 3 题，她应答 water, warmth, oxygen，三个都要说出来。\par
念第 4、5 题：答案 carbon dioxide (二氧化碳) 和 xylem (木质部，运水的管道)。\par
说：“答错的题我们待会儿重点讲。”\par
Say it:\par
pollination = pol-ih-NAY-shun\par
fertilisation = fer-tih-ly-ZAY-shun\par
ovule = OV-yool\par
germinate = JER-mih-nayt\par
xylem = ZY-lum\par
phloem = FLOH-um\par
stamen = STAY-mun}
\end{frame}

\begin{frame}{The one idea}
\relax
\vspace{6mm}
\FigGreenIdea\par\vspace{8mm}
Every Green Machines question asks about one of these three jobs.\par\vspace{2mm}
Say which job it is, then use that job's words.
\fref{Green Machines p.2}
\zh{%
屏幕：The one idea（一个核心思想）。三个方框：MAKE food 在叶子里制造食物 → MOVE materials 在植物体内运输物质 → MAKE new plants 用花产生新的植物。下面两句：每道 Green Machines 的题都在问这三件事中的一件；先说是哪一件，再用那一件的专用词。\par
背景：制造食物 = photosynthesis (光合作用)；运输 = water 走 xylem，glucose 走 phloem (韧皮部，运糖的管道)；繁殖 = 花、种子、萌发。}
\note{%
Teaching line, not a question slide.\par
时间 0:02–0:03。\par
MAKE food = photosynthesis (叶子) ；MOVE materials = transport：water 走 xylem，glucose 走 phloem；MAKE new plants = reproduction：flower、seed、germination。\par
问她：“testing a leaf for starch 属于哪一个？” 答：MAKE food，因为 starch 是储存起来的 glucose。\par
考试时先判断是哪一件事，再用那一块的术语，答案就不会跑题。\par
Script:\par
说：“整个单元其实只有三件事：make food、move materials、make new plants。”\par
问：“用碘液测叶子里有没有 starch (淀粉)，属于哪一件事？” 她应答：make food，因为 starch 是存起来的 glucose。\par
说：“考试拿到一道题，先想它属于哪一件，再用那一块的英文词来写，答案就不会跑题。”}
\end{frame}

% ---------------------------------------------------------- Loop 1
\begin{frame}{Quick review: Parts of a flower}
\relax
\twocol{0.42}{\vspace{1mm}\FigFlower[50mm]}{0.52}{%
\qline{1} Name parts A to F.
\qline{2} Which two parts make up the stamen?
\qline{3} Which three parts make up the carpel?
\qline{4} Which part becomes the fruit?}
\fref{Green Machines p.6}
\zh{%
屏幕：Quick review: Parts of a flower（快速回忆：花的结构）。左边是花的剖面图，标了字母 A–F。1. 说出 A 到 F 的名称。2. 哪两部分组成 stamen (雄蕊)？3. 哪三部分组成 carpel (心皮，也叫 pistil 雌蕊)？4. 哪一部分以后变成 fruit (果实)？\par
背景：A stigma 柱头，B style 花柱，C ovary 子房，D anther 花药，E filament 花丝，F sepal 萼片。雄性部分 stamen = anther + filament；雌性部分 carpel = stigma + style + ovary。受精后 ovary 变成果实，里面的 ovule 变成 seed (种子)。}
\note{%
1. A stigma, B style, C ovary, D anther, E filament, F sepal\par
2. anther and filament (D and E)\par
3. stigma, style and ovary (A, B and C)\par
4. the ovary (C)\par
Short solution:\par
4. fertilisation 之后 ovary wall 变成 fruit，里面的 ovule 变成 seed\par
Watch for:\par
1. sepal 和 petal 不要混：sepal 在花的最外面最下面，花还是花苞时保护它\par
3. carpel 也叫 pistil，两个名字都可以\par
Script:\par
说：“看图，从 A 到 F 说出名字，用英文。”\par
她应答：A stigma, B style, C ovary, D anther, E filament, F sepal。sepal 读 SEE-pul。\par
问：“Which two parts make the stamen?” 她答 anther and filament。再问：“Which three make the carpel?” 她答 stigma, style, ovary。\par
问：“Which part becomes the fruit?” 答 ovary。补一句：“The ovule becomes the seed.”\par
Say it:\par
sepal = SEE-pul\par
carpel = KAR-pul}
\end{frame}

\begin{frame}{Two ways to make new plants}
\relax
\renewcommand\arraystretch{1.35}
\begin{tabular}{@{}l l l@{}}
\toprule
 & \textbf{Asexual} & \textbf{Sexual}\\
\midrule
Parents & \gap & two\\
Gametes? & \gap & pollen + ovule\\
Offspring & \gap & \gap\\
Examples & runners, bulbs, cuttings & \gap\\
\bottomrule
\end{tabular}\par\vspace{4mm}
A spider plant grows plantlets on long stems.\par
Which type of reproduction is this? How do you know?
\fref{Green Machines p.3--5}
\zh{%
屏幕：Two ways to make new plants（植物繁殖的两种方式）。表格比较 Asexual (无性生殖) 和 Sexual (有性生殖)：几个 parents (亲本)？有没有 gametes (配子，即 pollen 和 ovule)？后代是否一样？例子：runners (匍匐茎)、bulbs (鳞茎)、cuttings (扦插)。下面问：spider plant (吊兰) 在长茎上长出 plantlets (小植株)，这是哪种生殖？你怎么知道？\par
背景：无性生殖只有一个亲本，没有配子，后代和亲本完全一样，叫 clones (克隆)。有性生殖需要两个亲本，花粉和胚珠结合，后代有 variation (变异，彼此不同)。}
\note{%
Teaching line, not a question slide.\par
时间 0:05–0:07。\par
空格答案：Asexual 一栏 one parent / no gametes / identical (clones)；Sexual 一栏 offspring show variation / flowers and seeds。\par
问题答案：asexual reproduction，因为只有 one parent，没有 gametes，plantlets 和 parent plant identical (clones)。\par
她在 booklet p.4 把 “This specific process involves gametes” 也打了勾，这是 \Lref{asexual}：asexual 一定没有 gametes。\par
Script:\par
说：“表格里有空格，你直接说出来。”\par
指 Asexual 一列：她应答 one parent、no gametes、offspring identical (clones)。指 Sexual 一列：offspring show variation、examples are flowers and seeds。\par
念下面的问题：“A spider plant grows plantlets. Which type? How do you know?” 她应答：asexual，因为 one parent、no gametes、plantlets are identical to the parent。\par
说：“你 booklet 第 4 页这里打错了一个勾：asexual 一定没有 gametes。”\par
Say it:\par
asexual = ay-SEK-shoo-ul\par
gametes = GAM-eets}
\end{frame}

\begin{frame}{Why variation matters: write it in this order}
\relax
\qline{1} The environment changes, e.g. a new disease arrives.
\qline{2} Offspring from sexual reproduction are \gap[22mm] from each other.
\qline{3} Some offspring have features that let them \gap[22mm].
\qline{4} These survive and \gap[22mm], so the species \gap[22mm].
\par\vspace{3mm}
{\bfseries A}: lines 1--2 \quad {\bfseries M}: links 2 to 3 \quad {\bfseries E}: the whole chain + the asexual contrast\par\vspace{3mm}
Why could a whole field of identical plants die from one disease?
\fref{Green Machines p.5}
\zh{%
屏幕：Why variation matters: write it in this order（为什么变异很重要：按这个顺序写）。1. 环境改变，比如出现一种新疾病。2. 有性生殖的后代彼此\_\_\_\_。3. 有些后代有能让它们\_\_\_\_的特征。4. 这些后代存活并\_\_\_\_，所以这个物种\_\_\_\_。底下：A 写到第 1–2 点，M 把第 2 点和第 3 点连起来，E 写完整条链并和无性生殖对比。最后问：为什么一整片完全一样的植物会被一种病全部杀死？\par
背景：这是 Explain 长答题的标准写法：环境变化 → 后代不同 → 有些能活 → 繁殖 → 物种延续。}
\note{%
Teaching line, not a question slide.\par
时间 0:07–0:08。\par
空格：different (show variation) / survive the change (resist the disease) / reproduce / continues。\par
问题答案：the plants are clones with the same genetic material，所以如果一株会得这种病，所有植株都会得，没有一株有抵抗力。\par
这是 Explain 长答题的标准链条。她 booklet p.7 写 “helps the species to live”，只到 A。要拿 E 必须写出 “some survive → reproduce → species continues”，并且和 asexual 对比。\par
Script:\par
说：“这一页是 Explain 题的写法，按 1、2、3、4 的顺序填空。”\par
她应答：different from each other；survive the change；reproduce；continues。\par
问：“Why could a whole field of identical plants die from one disease?” 她应答：they are clones，所以如果一株会得这种病，所有植株都会得。\par
说：“考试写这种题，四步都要写，只写 ‘helps the species survive’ 只能拿 A。”}
\end{frame}

\begin{frame}{From pollen to seed}
\relax
Put the cards in order.\par\vspace{2mm}
\qline{A} The pollen tube grows down the style.
\qline{B} Pollen lands on the stigma.
\qline{C} The ovule becomes a seed; the ovary becomes a fruit.
\qline{D} The pollen nucleus fuses with the ovule nucleus.
\qline{E} The pollen tube enters the ovule through the micropyle.
\qline{F} An insect carries pollen away from an anther.
\par\vspace{2mm}
Which cards are pollination? Which card is fertilisation?
\fref{Green Machines p.7, p.10}
\zh{%
屏幕：From pollen to seed（从花粉到种子）。把 A–F 六张卡片排顺序：A 花粉管沿着 style 往下长；B 花粉落到 stigma 上；C ovule 变成种子，ovary 变成果实；D 花粉的 nucleus 和 ovule 的 nucleus 融合；E 花粉管通过 micropyle (珠孔) 进入 ovule；F 昆虫从 anther 带走花粉。问：哪几张是 pollination？哪一张是 fertilisation？\par
背景：正确顺序 F, B, A, E, D, C。pollination 是 F 和 B（花粉从花药到柱头）。fertilisation 是 D（两个细胞核融合，形成 zygote 合子，以后长成 embryo 胚）。}
\note{%
Teaching line, not a question slide.\par
时间 0:08–0:10。\par
顺序：F, B, A, E, D, C。\par
Pollination = F 和 B (pollen 从 anther 到 stigma) 。Fertilisation = D (两个 nucleus 融合，形成 zygote) 。\par
提醒 \Lref{pollfert}：pollen 落在 stigma 上还不是 fertilisation。\par
D 之后 zygote 长成 embryo，在 seed 里面。\par
Script:\par
说：“六张卡片，按发生的先后排一下。”\par
她应答：F, B, A, E, D, C。\par
问：“Which cards are pollination?” 答 F and B。问：“Which card is fertilisation?” 答 D。\par
说：“记住：pollen 落在 stigma 上只是 pollination，还没有 fertilisation。”\par
Say it:\par
micropyle = MY-kroh-pile\par
zygote = ZY-goht\par
embryo = EM-bree-oh}
\end{frame}

\begin{frame}{Wind or insect pollination?}
\relax
\renewcommand\arraystretch{1.35}
\begin{tabular}{@{}l l l@{}}
\toprule
\textbf{Part} & \textbf{Wind-pollinated} & \textbf{Insect-pollinated}\\
\midrule
Petals & small, dull, no scent & \gap[30mm]\\
Stigma & \gap[30mm] & sticky, inside the flower\\
Anthers & hang outside the flower & \gap[30mm]\\
Pollen & \gap[30mm] & larger, sticky or spiky\\
\bottomrule
\end{tabular}\par\vspace{4mm}
Full marks: the feature \textbf{and} how it helps.\par\vspace{2mm}
Why do the anthers of grass hang outside the flower?
\fref{Green Machines p.8--9}
\zh{%
屏幕：Wind or insect pollination?（风媒还是虫媒授粉？）表格比较风媒花和虫媒花的 petals (花瓣)、stigma、anthers、pollen 四个部分，有空格要填。下面：满分要写出“特征”和“这个特征怎么帮助授粉”。问：为什么草的 anthers 挂在花的外面？\par
背景：风媒花（如草）：花瓣小、不鲜艳、没香味；柱头大、羽毛状、伸在外面接花粉；花药挂在外面让风吹走花粉；花粉小、轻、量多。虫媒花（如百合）：花瓣大、鲜艳、有香味和 nectar (花蜜)；柱头黏、在花里面；花药固定在里面，昆虫经过时蹭到；花粉较大、黏或带刺，能粘在昆虫身上。}
\note{%
Teaching line, not a question slide.\par
时间 0:10–0:12。\par
空格：large, brightly coloured, scented, with nectar / large and feathery, hanging outside / firmly attached inside the flower / small, light, made in huge amounts。\par
问题答案：so the wind can blow the pollen away from the anthers。只写 “they hang outside” 是 A；写出 “so the wind can carry the pollen” 才是 M。\par
这是 \Lref{adapt}：每个 feature 后面都要加 “so that ...”。\par
Script:\par
说：“把四个空格填上，说英文。”\par
她应答：insect petals large, brightly coloured, scented, with nectar；wind stigma large and feathery, hanging outside；insect anthers firmly inside the flower；wind pollen small, light, lots of it。\par
问：“Why do grass anthers hang outside the flower?” 她要答：so the wind can blow the pollen away。\par
说：“每个特征后面都要加 ‘so that…’，不然只有 A。”}
\end{frame}

\begin{frame}{Spot the error (1)}
\relax
A strawberry plant makes new plants from runners and from its flowers. Compare the two.\par\vspace{2mm}
\samcard{146mm}{\flawed{Runners: the new plant gets gametes from the parent, so it is a bit different from it. Flowers: pollen is carried from the stigma to the anther. This is called fertilisation. The seeds show variation because two parents are involved.}}\par\vspace{2mm}
\textbf{Sam's answer has 3 errors.} Find each one, fix it, and say why it loses marks.
\fref{Green Machines p.3--7}
\zh{%
屏幕：Spot the error (1)（找错误 1）。题目：草莓用 runners 和花两种方式产生新植株，比较这两种方式。Sam 的答案（手写体）：Runners：新植株从亲本得到 gametes，所以和亲本有点不同。Flowers：花粉从 stigma 被带到 anther，这叫 fertilisation。种子有变异，因为有两个亲本。下面：Sam 的答案有 3 个错误，找出来、改正、并说为什么会丢分。\par
背景：错误 1：runners 是无性生殖，没有 gametes，新植株和亲本完全一样。错误 2：花粉方向反了，应该从 anther 到 stigma。错误 3：这个过程叫 pollination，不叫 fertilisation。最后一句是对的。}
\note{%
1. runners have no gametes: the new plant is identical to the parent, a clone (\Lref{asexual})\par
2. pollen is carried from the anther to the stigma, not the other way (\Lref{pollfert})\par
3. that transfer is pollination, not fertilisation (\Lref{pollfert})\par
Short solution:\par
Runners are asexual: one parent, no gametes, so the new plants are identical to the parent. Flowers: pollen is carried from the anther to the stigma (pollination); the pollen nucleus fuses with the ovule nucleus (fertilisation). The seeds show variation because two parents are involved.\par
Watch for:\par
最后一句是对的。让她先判断每一句，而不是看到 Sam 就全部改掉。\par
Script:\par
说：“Sam 是一个虚构的同学。他的答案里有 3 个错误，你当老师，一句一句检查。”\par
等她找。她应找出：runners have no gametes and make identical clones；pollen goes from anther to stigma；that is pollination, not fertilisation。\par
问：“最后一句对不对？” 她应答：对，两个亲本所以有变异。\par
说：“不要看到 Sam 就全部改，要一句一句判断。”}
\end{frame}

\begin{frame}{Quick review: Seeds}
\relax
\twocol{0.50}{\vspace{2mm}\FigSeed[66mm]}{0.46}{%
\qline{1} Name parts P to S.
\qline{2} Name the three conditions for germination.
\qline{3} Why are seeds dispersed far from the parent plant?}
\fref{Green Machines p.11--13}
\zh{%
屏幕：Quick review: Seeds（快速回忆：种子）。左边是剖开的豆子，标了 P、Q、R、S。1. 说出 P 到 S 的名称。2. 说出 germination (萌发) 的三个条件。3. 为什么种子要被 dispersed (传播) 到离母株远的地方？\par
背景：P testa 种皮（保护种子的硬外壳，吸水后变软），Q plumule 胚芽（向上长成茎叶），R radicle 胚根（向下长成根），S cotyledon 子叶（储存养分，在叶子能光合作用前提供能量）。三个条件：water, warmth, oxygen。传播到远处是为了不和母株竞争 light、water、minerals (矿物质) 和 space。}
\note{%
1. P testa, Q plumule, R radicle, S cotyledon\par
2. water, warmth, oxygen\par
3. so the new plant does not compete with the parent for light, water, minerals and space\par
Short solution:\par
1. testa 是保护种子的硬外皮，water 让它变软；radicle 向下长成 root，固定植物并吸水；plumule 向上长成 shoot；cotyledon 是 food store，在叶子能 photosynthesise 之前提供能量\par
2. oxygen 用于 respiration，释放生长需要的能量\par
Watch for:\par
1. booklet 印的是 “Radical”，考试要拼 radicle\par
3. 只写 “to spread out” 是 A；要写出 competition 和资源名称\par
Script:\par
说：“看豆子的图，P 到 S 分别是什么？”\par
她应答：P testa, Q plumule, R radicle, S cotyledon。提醒：booklet 印成 Radical，考试要拼 radicle。\par
问第 2 题：她答 water, warmth, oxygen。\par
问第 3 题：她要说出 competition，以及 light, water, minerals, space。只说 ‘spread out’ 不够。\par
Say it:\par
plumule = PLOO-myool\par
radicle = RAD-ih-kul\par
cotyledon = kot-ih-LEE-dun\par
dispersed = dis-PERST}
\end{frame}

\begin{frame}{Practice}
\relax
\qline{Q1} A lily has large white petals, a strong scent and nectar. Is it wind- or insect-pollinated? Give two reasons.\slv{A}
\qline{Q2} A grower takes cuttings from her best rose bush instead of planting its seeds. Explain why.\slv{M}
\qline{Q3} Cress seeds are put in four tubes at 20\,\textdegree C: A on dry cotton wool; B on wet cotton wool; C on wet cotton wool in a fridge; D under boiled water with a layer of oil on top. Which tubes germinate? Explain the result for D.\slv{E}
\fref{Green Machines p.5, p.8, p.12}
\zh{%
屏幕：Practice（练习）。Q1 [A] 百合花有大的白色花瓣、浓香味和花蜜，是风媒还是虫媒？给两个理由。Q2 [M] 种花的人从最好的玫瑰上剪枝扦插，而不是种它的种子，为什么？Q3 [E] 水芹种子放在四支试管里（20 °C）：A 干棉花；B 湿棉花；C 湿棉花放冰箱；D 泡在煮过的水里，上面盖一层油。哪几管会萌发？解释 D 的结果。\par
背景：Q2：扦插是无性生殖，新植株和母株一样，保留好的特征；种子来自有性生殖，有变异。Q3：只有 B 萌发；A 缺水，C 太冷，D 缺氧——水煮过后溶解的氧气被赶走，油层阻止空气再溶进去；种子需要氧气进行 respiration (呼吸作用) 释放能量。}
\note{%
Q1 insect-pollinated: large white petals and scent attract insects; nectar is food for insects\par
Q2 cuttings are asexual, so the new plants are identical to the best bush and keep its features; seeds come from sexual reproduction and show variation\par
Q3 only B germinates; D has no oxygen, so the seed cannot respire to release energy for growth\par
Short solution:\par
Q3 A 没有 water；C 没有 warmth (fridge 里温度太低) ；D：boiling 把水里溶解的 oxygen 赶走，oil 层不让空气重新溶进去\par
Watch for:\par
Q1 白色花瓣也算 “large, showy”，不一定要彩色\par
Q3 C 的标准答案写 “too cold / no warmth”。D 一定要写 oxygen 的用途 (respiration 释放能量) ，否则只有 M\par
Q3 这个 boiled water + oil 的设计和 Lesson 2 的 rusty nail 完全一样，等一下会用到\par
Script:\par
说：“三道练习题，[A] [M] [E] 是难度。先自己说答案，再看我补充。”\par
Q1 她应答：insect-pollinated，因为 large white petals and scent attract insects，nectar is food for insects。\par
Q2 她应答：cuttings are asexual, so new plants are identical and keep the good features; seeds show variation。\par
Q3 她应答：only B；D has no oxygen, so the seed cannot respire to release energy。如果她只说 ‘no oxygen’，追问：“oxygen 是用来做什么的？”\par
说：“记住 D 这个设计 (煮过的水 + 油层 = 没有氧气) ，第二课生锈实验会再用到。”}
\end{frame}

% ---------------------------------------------------------- Loop 2
\begin{frame}{Quick review: Leaf structure}
\relax
\hspace*{9mm}\FigLeaf[70mm]\hfill
\begin{minipage}[t]{50mm}\vspace{0pt}\raggedright
\qline{1} Name A to H.
\qline{2} Which layer does most of the photosynthesis?
\end{minipage}
\fref{Green Machines p.15}
\zh{%
屏幕：Quick review: Leaf structure（快速回忆：叶片结构）。叶片横切面图，标了 A–H。1. 说出 A 到 H。2. 哪一层进行最多的 photosynthesis？\par
背景：A waxy cuticle 蜡质角质层（防止水分流失）；B upper epidermis 上表皮；C palisade layer 栅栏层（细胞又高又密，含最多 chloroplasts 叶绿体）；D spongy layer 海绵层（细胞之间有空隙，让气体流动）；E vein 叶脉，里面有 xylem 和 phloem；F lower epidermis 下表皮；G guard cell 保卫细胞；H stoma 气孔（复数 stomata）。}
\note{%
1. A waxy cuticle, B upper epidermis, C palisade layer, D spongy layer, E vein (xylem and phloem), F lower epidermis, G guard cell, H stoma\par
2. the palisade layer (C)\par
Short solution:\par
2. palisade cells 又高又窄、排得很紧，在叶子上部，含最多 chloroplasts\par
Watch for:\par
1. H 的复数是 stomata；拼写常错\par
1. 她 booklet p.15 把 “transports water and minerals” 填成 spongy layer。应该是 vein 里面的 xylem (\Lref{xylem})\par
Script:\par
说：“叶子的横切面，从上到下说出 A 到 H。”\par
她应答：A waxy cuticle, B upper epidermis, C palisade layer, D spongy layer, E vein, F lower epidermis, G guard cell, H stoma。\par
问：“Which layer does most photosynthesis?” 答 palisade layer，因为它在上面、chloroplasts 最多。\par
说：“vein 里面有 xylem 和 phloem。问 ‘哪个运水’ 要答 xylem，只写 vein 不够。”\par
Say it:\par
cuticle = KYOO-tih-kul\par
epidermis = ep-ih-DER-mis\par
palisade = pal-ih-SAYD\par
stomata = STOH-muh-tuh\par
chloroplasts = KLOR-uh-plasts}
\end{frame}

\begin{frame}{Photosynthesis: the word equation}
\relax
\vspace{3mm}
\hbox to\textwidth{\hfil\gap[30mm]\ + \gap[22mm]\ \ $\xrightarrow[\rule{20mm}{0.4pt}]{\rule{24mm}{0.4pt}}$\ \ \gap[22mm]\ + \gap[22mm]\hfil}\par\vspace{6mm}
Where does each starting substance enter the plant?\par\vspace{2mm}
\qline{} carbon dioxide: through the \gap[26mm] in the leaf
\qline{} water: root hairs $\to$ \gap[22mm] $\to$ leaf\par\vspace{3mm}
What does chlorophyll do in this reaction?
\fref{Green Machines p.20}
\zh{%
屏幕：Photosynthesis: the word equation（光合作用文字方程式）。第一行是方程式的空格：\_\_\_ + \_\_\_ →（箭头上下各一个空）\_\_\_ + \_\_\_。下面问：每种原料从哪里进入植物？carbon dioxide 通过叶子上的\_\_\_\_；water 从 root hairs (根毛) → \_\_\_\_ → 叶子。最后问：chlorophyll (叶绿素) 在这个反应里起什么作用？\par
背景：carbon dioxide + water → glucose + oxygen，箭头上写 light energy (光能)，下面写 chlorophyll。二氧化碳从 stomata 进入；水靠 osmosis (渗透作用) 进入根毛，再经 xylem 上到叶子。chlorophyll 吸收光能（booklet 说它“加快反应”，两种说法都接受）。}
\note{%
Teaching line, not a question slide.\par
时间 0:19–0:21。\par
方程式：carbon dioxide + water → glucose + oxygen；箭头上面写 light energy，下面写 chlorophyll。\par
carbon dioxide 从 stomata 进入；water 从 root hairs 进入 (osmosis)，经 xylem 到叶子。\par
chlorophyll：absorbs light energy (真实的科学) 。booklet p.20 写的是 chlorophyll “speeds up the reaction”，按学校的说法两种都可以接受。\par
\Lref{photoeq}：light 和 chlorophyll 不是反应物，写在箭头上下。\par
Script:\par
说：“这是光合作用的文字方程式，你把四个空和箭头上下的两个空说出来。”\par
她应答：carbon dioxide + water → glucose + oxygen；arrow 上面 light energy，下面 chlorophyll。\par
问：“Where does carbon dioxide get in?” 答 through the stomata。问水的路线：root hairs → xylem → leaf。\par
问：“What does chlorophyll do?” 答 absorbs light energy。\par
说：“light 和 chlorophyll 不是原料，写在箭头上下，不要写进加号里。”\par
Say it:\par
chlorophyll = KLOR-uh-fil\par
osmosis = oz-MOH-sis}
\end{frame}

\begin{frame}{Respiration: the same substances, reversed}
\relax
\vspace{2mm}
\hbox to\textwidth{\hfil glucose + \gap[22mm]\ \ $\longrightarrow$\ \ \gap[30mm]\ + \gap[22mm]\ \ (+ energy)\hfil}\par\vspace{5mm}
\renewcommand\arraystretch{1.35}
\hbox to\textwidth{\hfil\begin{tabular}{@{}l l l@{}}
\toprule
 & \textbf{Photosynthesis} & \textbf{Respiration}\\
\midrule
Day & \gap[24mm] & \gap[24mm]\\
Night & \gap[24mm] & \gap[24mm]\\
\bottomrule
\end{tabular}\hfil}\par\vspace{4mm}
Which process happens in every living cell, all the time?
\fref{Green Machines p.23}
\zh{%
屏幕：Respiration: the same substances, reversed（呼吸作用：同样的物质，反过来）。方程式空格：glucose + \_\_\_ → \_\_\_ + \_\_\_ (+ energy)。下面表格：白天和晚上，photosynthesis 和 respiration 是否进行（active / not active）。问：哪个过程在每个活细胞里、任何时候都在进行？\par
背景：glucose + oxygen → carbon dioxide + water (+ energy)，在 mitochondria (线粒体) 里进行。白天两个都进行；晚上没有光，只有 respiration。所以植物也是一直在呼吸。}
\note{%
Teaching line, not a question slide.\par
时间 0:21–0:23。\par
方程式：glucose + oxygen → carbon dioxide + water (+ energy)。在 mitochondria 里进行。\par
表格：Day 两个都 active；Night photosynthesis not active，respiration active。\par
问题答案：respiration。\par
\Lref{respire}：植物白天也 respire，只是白天 photosynthesis 更快，所以白天净放出 oxygen。\par
Script:\par
说：“呼吸作用正好是光合作用倒过来。”\par
她应答：glucose + oxygen → carbon dioxide + water + energy。\par
指表格：Day 两个都 active；Night photosynthesis not active, respiration active。\par
问：“Which process happens in every living cell, all the time?” 答 respiration。\par
说：“很多学生写 ‘plants respire at night’，错，植物白天晚上都 respire。”\par
Say it:\par
mitochondria = my-toh-KON-dree-uh}
\end{frame}

\begin{frame}{Testing a leaf for starch}
\relax
Give the reason for each step.\par\vspace{2mm}
\qline{1} Put the leaf in boiling water for 1 minute.
\qline{2} Put it in hot ethanol, standing in a beaker of hot water.
\qline{3} Rinse it in water.
\qline{4} Add iodine solution. Blue-black means \gap[34mm].
\par\vspace{3mm}
Why is the ethanol heated in hot water, not over a Bunsen flame?
\fref{Green Machines p.20}
\zh{%
屏幕：Testing a leaf for starch（测叶子里的淀粉）。给每一步写出原因：1. 把叶子放进沸水 1 分钟。2. 放进热的 ethanol (乙醇)，乙醇管放在一杯热水里。3. 用水冲洗。4. 滴 iodine (碘液)，变成蓝黑色说明\_\_\_\_。最后问：为什么乙醇用热水加热，而不是放在本生灯火焰上？\par
背景：1 杀死叶片、停止反应；2 去掉叶绿素，叶子变白，才能看清颜色变化；3 让变脆的叶子变软、洗掉乙醇；4 蓝黑色 = 有 starch；没有淀粉时保持棕黄色。乙醇易燃，不能靠近明火。}
\note{%
Teaching line, not a question slide.\par
时间 0:23–0:25。\par
1. kills the leaf and stops its reactions (so the starch stays)\par
2. removes the chlorophyll, so the colour change can be seen\par
3. softens the leaf and washes off the ethanol\par
4. starch is present；没有 starch 时保持 brown / yellow-brown\par
问题答案：ethanol is flammable，不能靠近明火。\par
实验题常考 “why each step”，她最容易只写步骤不写原因。\par
Script:\par
说：“实验题常考 ‘每一步为什么做’，这一页四步，每步说一个原因。”\par
她应答：1 kills the leaf；2 removes the chlorophyll so the colour change can be seen；3 softens the leaf；4 starch is present。\par
问最后一题：答 ethanol is flammable (易燃)。\par
说：“只写步骤不写原因，实验题拿不到 M。”\par
Say it:\par
ethanol = ETH-uh-nol\par
iodine = EYE-uh-dyne}
\end{frame}

\begin{frame}{Workbook: the foil leaf}
\relax
\twocol{0.40}{\vspace{3mm}\FigVarLeaf[56mm]}{0.56}{%
A variegated leaf: J green, K white. Foil L covered part of it for 2 days in the light. Then the leaf was tested with iodine.\par\vspace{3mm}
Which areas turn blue-black? Explain each area.}
\fref{Green Machines p.22}
\zh{%
屏幕：Workbook: the foil leaf（练习册题：盖锡箔的叶子）。一片 variegated leaf (斑叶，有绿有白)：J 是绿色部分，K 是白色部分，L 是盖了锡箔的部分，在光下放 2 天后用碘液测试。问：哪些部分变蓝黑？逐个解释。\par
背景：J 蓝黑：有 chlorophyll 又有光，能光合作用，把 glucose 存成 starch。K 保持棕色：没有 chlorophyll。L 保持棕色：没有光（即使是绿的）。这是她 booklet 第 22 页没做的题。}
\note{%
(a) J (green, not covered): blue-black\par
(b) K (white): stays brown\par
(c) L (under the foil): stays brown, even where it is green\par
Short solution:\par
J 有 chlorophyll 也有 light，进行 photosynthesis，glucose 变成 starch 储存\par
K 没有 chlorophyll，不能 photosynthesise，没有 starch\par
L 没有 light，不能 photosynthesise，没有 starch\par
Watch for:\par
这是 booklet p.22 第 2 题，她没写。每个区域都要说 “缺什么”：K 缺 chlorophyll，L 缺 light\par
Script:\par
说：“这是你 booklet 第 22 页没做的题，我们现在做。”\par
问：“Which areas turn blue-black?” 她应答：only J。\par
让她解释：K has no chlorophyll；L has no light；所以都不能 photosynthesise，没有 starch。\par
说：“每个区域都要说清楚缺了什么：K 缺 chlorophyll，L 缺 light。”\par
Say it:\par
variegated = VAIR-ee-uh-gay-tid}
\end{frame}

\begin{frame}{Spot the error (2)}
\relax
Write the word equation for photosynthesis. Then say how you would show that a leaf has made food.\par\vspace{2mm}
\samcard{146mm}{\flawed{oxygen + water makes glucose + carbon dioxide (light above the arrow). Boil the leaf, put it in hot ethanol, rinse it, then add iodine. If it goes blue-black, the leaf has made glucose. Plants photosynthesise in the day and respire at night.}}\par\vspace{2mm}
\textbf{Sam's answer has 3 errors.} Find each one, fix it, and say why it loses marks.
\fref{Green Machines p.20--23}
\zh{%
屏幕：Spot the error (2)。题目：写出光合作用的文字方程式，并说明怎么证明叶子制造了食物。Sam：氧气 + 水 生成 葡萄糖 + 二氧化碳（箭头上写光）。把叶子煮一下，放进热乙醇，冲洗，滴碘液；如果变蓝黑，说明叶子制造了 glucose。植物白天光合作用，晚上呼吸。3 个错误。\par
背景：错误 1：氧气和二氧化碳写反了。错误 2：碘液测的是 starch，不是 glucose。错误 3：植物白天晚上都呼吸。实验步骤那句是对的。}
\note{%
1. carbon dioxide + water → glucose + oxygen: Sam swapped oxygen and carbon dioxide (\Lref{photoeq})\par
2. blue-black shows starch, not glucose (\Lref{iodine})\par
3. plants respire all the time, day and night (\Lref{respire})\par
Short solution:\par
carbon dioxide + water → glucose + oxygen, with light energy above the arrow and chlorophyll below. Blue-black with iodine shows the leaf contains starch. Plants respire all the time; they photosynthesise only in the light.\par
Watch for:\par
实验步骤那一句是对的。\par
Script:\par
说：“又是 Sam，3 个错误。”\par
她应找出：carbon dioxide and oxygen are swapped；iodine shows starch, not glucose；plants respire all the time。\par
问：“实验步骤那句对吗？” 答：对。}
\end{frame}

\begin{frame}{Practice}
\relax
\qline{Q1} Two potted plants are kept in the light for 2 days, each under a sealed bell jar. Jar 1 also holds a dish of soda lime, which absorbs carbon dioxide. Jar 2 holds a dish of water. A leaf from each plant is tested with iodine. Predict and explain the results.\slv{M}
\qline{Q2} In a greenhouse full of plants, the carbon dioxide level rises at night and falls during the day. Explain why.\slv{E}
\fref{Green Machines p.21, p.23}
\zh{%
屏幕：Practice。Q1 [M] 两盆植物在光下放 2 天，各罩在密封的钟罩里。1 号罩里还放了一碟 soda lime (碱石灰，会吸收二氧化碳)；2 号罩里放一碟水。各取一片叶子测碘，预测并解释结果。Q2 [E] 种满植物的温室里，二氧化碳浓度晚上升高、白天下降，为什么？\par
背景：Q1：1 号没有二氧化碳，不能光合作用，叶子保持棕色；2 号是对照，变蓝黑。Q2：晚上只有呼吸，放出二氧化碳；白天光合作用吸收二氧化碳的速度比呼吸放出的快。}
\note{%
Q1 plant 1 stays brown (no starch); plant 2 turns blue-black (starch)\par
Q2 at night only respiration happens, which gives out carbon dioxide; in the day photosynthesis takes in carbon dioxide faster than respiration gives it out\par
Short solution:\par
Q1 jar 1 没有 carbon dioxide，不能 photosynthesise，不能做 glucose，所以没有 starch；jar 2 有 carbon dioxide，是 control\par
Q2 要写清两个过程，以及白天 “photosynthesis is faster than respiration”\par
Watch for:\par
Q1 她 booklet p.21 只写了 “there won't be photosynthesis”，没写原因 (缺 carbon dioxide)，只能拿 A\par
Q2 写 “plants respire at night” 只答了一半；只写 “plants breathe out carbon dioxide” 不给分\par
Script:\par
念 Q1，问：“Predict and explain.” 她应答：plant 1 stays brown because there is no carbon dioxide, so no photosynthesis, no starch；plant 2 turns blue-black。\par
追问：“Jar 2 有什么用？” 答：it is the control (对照)。\par
念 Q2，她应答：at night only respiration, which gives out carbon dioxide；in the day photosynthesis uses carbon dioxide faster than respiration makes it。\par
说：“只写 ‘plants respire at night’ 只答了一半。”}
\end{frame}

% ---------------------------------------------------------- Loop 3
\begin{frame}{Quick review: Transport words}
\relax
Match each word to its job.\par\vspace{2mm}
\begin{tabular}{@{}p{38mm}p{108mm}@{}}
\textbf{1} xylem & \textbf{A} carries glucose from the leaves\\
\textbf{2} phloem & \textbf{B} lets gases and water vapour in and out of the leaf\\
\textbf{3} root hair & \textbf{C} absorbs water from the soil\\
\textbf{4} stomata & \textbf{D} carries water and minerals up the plant\\
\textbf{5} guard cells & \textbf{E} change shape to open and close the stomata\\
\end{tabular}
\fref{Green Machines p.16--17}
\zh{%
屏幕：Quick review: Transport words（快速回忆：运输相关的词）。把左边 1–5 和右边 A–E 配对：1 xylem 2 phloem 3 root hair 4 stomata 5 guard cells；A 把 glucose 从叶子运走 B 让气体和水蒸气进出叶子 C 从土壤吸水 D 把水和矿物质往上运 E 改变形状来打开和关闭 stomata。\par
背景：答案 1D 2A 3C 4B 5E。注意：booklet 第 16 页说 guard cells “膨胀时关闭气孔”，实际上是相反的（吸水膨胀时打开）。考试不太会考方向，如果考了按 booklet 写。}
\note{%
1. D\par
2. A\par
3. C\par
4. B\par
5. E\par
Short solution:\par
记法：xylem 运 water 向上；phloem 运 food (glucose) 到需要或储存的地方\par
Watch for:\par
2. phloem 的拼写：ph 开头，-oem 结尾\par
5. booklet p.16 说 guard cells “swell to close” stomata。真实情况是 guard cells 吸水 swell 时 stomata 打开。考试如果问 swell 和 open 的关系，按 booklet 写，并提醒她两种说法都见过\par
Script:\par
说：“配对，说数字加字母。”\par
她应答：1 D, 2 A, 3 C, 4 B, 5 E。\par
说：“记法：xylem 运 water 往上，phloem 运 food，也就是 glucose。”\par
提醒拼写：phloem 是 ph 开头。}
\end{frame}

\begin{frame}{How water travels through a plant}
\relax
\vspace{2mm}
soil water\par
\quad $\downarrow$ by \gap[26mm]\par
root hair cell\par
\quad $\downarrow$\par
\gap[26mm] in the root and up the stem\par
\quad $\downarrow$\par
leaf cells: water evaporates\par
\quad $\downarrow$\par
out through the \gap[26mm] as water vapour\par\vspace{3mm}
What is the loss of water vapour from the leaves called?
\fref{Green Machines p.16}
\zh{%
屏幕：How water travels through a plant（水在植物里怎么走）。一条从上往下的流程：土壤里的水 → 通过\_\_\_\_进入根毛细胞 → 在根里和茎里的\_\_\_\_ → 叶细胞：水蒸发 → 以水蒸气的形式从\_\_\_\_出去。问：叶子失去水蒸气叫什么？\par
背景：空格依次是 osmosis、xylem、stomata；叶子失去水蒸气叫 transpiration (蒸腾作用)。}
\note{%
Teaching line, not a question slide.\par
时间 0:31–0:33。\par
空格：osmosis / xylem / stomata。\par
问题答案：transpiration。\par
Script:\par
说：“跟着箭头走一遍，把空格说出来。”\par
她应答：osmosis；xylem；stomata。\par
问：“What is the loss of water vapour from the leaves called?” 答 transpiration。\par
Say it:\par
transpiration = tran-spih-RAY-shun}
\end{frame}

\begin{frame}{Why the water moves up: write it in this order}
\relax
A celery stalk stands in blue dye for 2 days. The leaves turn blue.\par\vspace{2mm}
\qline{1} Water \gap[26mm] from the leaves through the stomata.
\qline{2} This \gap[20mm] more water up the \gap[20mm].
\qline{3} The dye travels with the water, so the \gap[20mm] turn blue.
\qline{4} Cut the stalk: the blue dots show where the \gap[20mm] are.
\par\vspace{3mm}
{\bfseries A}: dye goes up \quad {\bfseries M}: names xylem \quad {\bfseries E}: transpiration causes the pull
\fref{Green Machines p.17}
\zh{%
屏幕：Why the water moves up: write it in this order（水为什么往上走：按顺序写）。芹菜放在蓝色色素水里两天，叶子变蓝。1. 水从叶子通过 stomata \_\_\_\_。2. 这会把更多水沿着\_\_\_\_往上\_\_\_\_。3. 色素跟着水走，所以\_\_\_\_变蓝。4. 切开茎，蓝点显示\_\_\_\_的位置。A：色素往上走；M：说出 xylem；E：说出是 transpiration 造成的拉力。\par
背景：答案 evaporates；pulls / xylem；leaves；xylem vessels。她 booklet 第 17 页只写了“水被运上去、叶子变色”，没有 xylem、没有原因，只有 A。}
\note{%
Teaching line, not a question slide.\par
时间 0:33–0:35。\par
空格：evaporates (is lost) / pulls / xylem / leaves / xylem vessels。\par
她 booklet p.17 的答案是 “The water in the cylinder are being transport up to the stem. The food colour will change the colour of plant”，只到 A：没有 xylem，没有原因。\par
\Lref{transp}：不能写 “the plant sucks / drinks”，要写 transpiration 造成 pull。\par
Script:\par
说：“芹菜实验，按四步解释。”\par
她应答：water evaporates from the leaves；this pulls more water up the xylem；the leaves turn blue；the dots show the xylem。\par
说：“千万不要写 ‘the plant sucks up water’ 或 ‘it is thirsty’，要写 transpiration 造成的 pull。”}
\end{frame}

\begin{frame}{Osmosis}
\relax
\twocol{0.32}{\vspace{1mm}\FigMembrane[42mm]}{0.64}{%
Left: sugar solution. Right: water.\par\vspace{2mm}
\qline{1} Which shape is a water molecule?
\qline{2} Which way do the water molecules move?
\qline{3} Why can the sugar molecules not cross?}
\par\vspace{3mm}
Osmosis is the movement of \gap[18mm] from a \gap[16mm] concentration of water to a \gap[16mm] concentration of water through a \gap[30mm] membrane.
\fref{Green Machines p.19}
\zh{%
屏幕：Osmosis（渗透作用）。左边是放大的半透膜图，左侧是糖溶液，右侧是水。1. 哪种形状是水分子？2. 水分子往哪边走？3. 为什么糖分子过不去？下面填空：Osmosis 是\_\_\_\_从\_\_\_\_浓度的水到\_\_\_\_浓度的水，通过\_\_\_\_膜的移动。\par
背景：小圆点是水分子，六边形是糖分子。水从右往左走（进入糖溶液）；糖分子太大，穿不过膜上的小孔。定义：water；high；low；semi-permeable (半透膜)。注意说的是“水的浓度”。}
\note{%
Teaching line, not a question slide.\par
时间 0:35–0:37。\par
1. the small circles (六边形是 sugar molecules) \par
2. from right to left, into the sugar solution\par
3. the sugar molecules are too large to fit through the membrane\par
定义空格：water / high / low / semi-permeable。\par
\Lref{osmosis}：说 “sugar moves” 或者漏掉 membrane 都会丢分；“high to low” 一定要说是 water 的浓度。\par
Script:\par
说：“看放大图，小圆圈和六边形哪个是水？”\par
她应答：small circles are water；water moves from right to left into the sugar solution；sugar molecules are too big to fit through。\par
让她念完整定义：Osmosis is the movement of water from a high concentration of water to a low concentration of water through a semi-permeable membrane。\par
说：“‘high to low’ 后面一定要说 ‘of water’，也一定要说 membrane。”\par
Say it:\par
semi-permeable = sem-ee-PER-mee-uh-bul\par
molecule = MOL-ih-kyool}
\end{frame}

\begin{frame}{Workbook: the Visking tubing model}
\relax
\twocol{0.16}{\vspace{1mm}\hspace*{4mm}\FigVisking[9mm]}{0.80}{%
The tubing Z sits in a beaker of water. Over 30 minutes the liquid rises up X.\par\vspace{2mm}
\qline{1} Which part acts like the root hair? Which acts like the stem?
\qline{2} Explain why the liquid rises.
\qline{3} A root hair is long and thin. How does this help it?}
\fref{Green Machines p.17--18}
\zh{%
屏幕：Workbook: the Visking tubing model（练习册题：透析袋模型）。图：X 毛细管，Y 加了色素的蔗糖溶液，Z 扎紧的 Visking tubing (透析袋)。透析袋放在一杯水里，30 分钟后液体沿 X 上升。1. 哪部分像根毛？哪部分像茎？2. 解释液体为什么上升。3. 根毛又长又细，有什么好处？\par
背景：Z 像根毛（半透膜），X 像茎。水通过渗透作用从烧杯（水浓度高）进入蔗糖溶液，体积变大把液体推上去。根毛长而细 → surface area (表面积) 大 → 吸水多。}
\note{%
1. Z (the Visking tubing) is the root hair; X (the capillary tube) is the stem\par
2. water moves into the sucrose solution by osmosis, from the beaker (high concentration of water) through the membrane, so the liquid is pushed up the tube\par
3. it has a large surface area, so it can absorb more water by osmosis\par
Short solution:\par
2. Visking tubing 是 semi-permeable membrane：water 可以过，sucrose 不能过\par
Watch for:\par
1. 这是 booklet p.18 的 Conclusions，她没写\par
3. 只写 “absorbs water” 不够，要说 surface area\par
Script:\par
说：“这是 booklet 第 18 页的老师演示，你当时没写结论，我们现在说。”\par
她应答：Z is the root hair, X is the stem；water moves into the sucrose solution by osmosis, so the liquid rises；long and thin gives a large surface area, so it absorbs more water。\par
说：“第 3 问一定要说 surface area。”\par
Say it:\par
sucrose = SOO-krohs}
\end{frame}

\begin{frame}{Spot the error (3)}
\relax
A white carnation stands in blue dye. A day later its petals are blue. Explain.\par\vspace{2mm}
\samcard{146mm}{\flawed{The blue water was carried up the phloem tubes. The flower sucks the water up because it is thirsty. Water always gets into a plant by osmosis, from a low concentration of water to a high concentration of water.}}\par\vspace{2mm}
\textbf{Sam's answer has 3 errors.} Find each one, fix it, and say why it loses marks.
\fref{Green Machines p.16--19}
\zh{%
屏幕：Spot the error (3)。题目：白色康乃馨插在蓝色色素水里，一天后花瓣变蓝，解释。Sam：蓝色的水被 phloem 运上去。花因为口渴把水吸上去。水总是靠渗透作用进入植物，从水浓度低的地方到水浓度高的地方。3 个错误。\par
背景：错误 1：应该是 xylem。错误 2：不是“吸、口渴”，是水从 stomata 蒸发（transpiration）把水拉上去。错误 3：渗透作用方向反了，应该从高到低。}
\note{%
1. the water goes up the xylem, not the phloem (\Lref{xylem})\par
2. not “sucks because it is thirsty”: water evaporates through the stomata (transpiration) and pulls water up the xylem (\Lref{transp})\par
3. osmosis is from a high to a low concentration of water (\Lref{osmosis})\par
Short solution:\par
Water evaporates from the flower and leaves through the stomata (transpiration). This pulls water, with the dye, up the xylem, so the petals turn blue. Water enters root hairs by osmosis, from a high to a low concentration of water.\par
Watch for:\par
第 3 句只算一个 error (方向反了) 。\par
Script:\par
说：“Sam 的第三份答案，3 个错误。”\par
她应找出：xylem not phloem；not sucking, transpiration pulls the water up；osmosis is from high to low concentration of water。}
\end{frame}

\begin{frame}{Practice}
\relax
\qline{Q1} Name the process that moves water into a root hair, and the tubes that carry it to the leaves.\slv{A}
\qline{Q2} In a Visking tubing model the liquid is 12\,mm up the tube at the start, 30\,mm after 10 minutes and 48\,mm after 20 minutes.\par
(a) Calculate the average rate of rise in mm per minute.\par
(b) Explain why the liquid rises.\slv{M}
\qline{Q3} A plant wilts when its roots are put in very salty water. Explain why.\slv{E}
\fref{Green Machines p.16--19}
\zh{%
屏幕：Practice。Q1 [A] 说出把水送进根毛的过程，以及把水运到叶子的管道。Q2 [M] 透析袋模型里，液体一开始在管里 12 mm 高，10 分钟后 30 mm，20 分钟后 48 mm。(a) 算平均上升速度（mm/分钟）。(b) 解释液体为什么上升。Q3 [E] 把植物的根放进很咸的水里，植物会萎蔫 (wilt)，解释。\par
背景：Q1 osmosis；xylem。Q2 (a) (48 − 12) ÷ 20 = 1.8 mm 每分钟；常见错误是 48 ÷ 20，忘了减去开始的 12。Q3 咸水的水浓度比根细胞低，所以水通过渗透作用从根里流出来，植物失水萎蔫。}
\note{%
Q1 osmosis; xylem\par
Q2 (a) 1.8 mm per minute\par
Q2 (b) water enters the sugar solution by osmosis through the membrane, so the volume rises\par
Q3 the salty water has a lower concentration of water than the root cells, so water moves out of the roots by osmosis; the plant loses water and wilts\par
Short solution:\par
Q2 (a) rise $= 48 - 12 = 36\ \mathrm{mm}$; rate $= \dfrac{36\ \mathrm{mm}}{20\ \mathrm{min}} = 1.8$ mm per minute\par
Q3 先比较 water concentration (外面低、里面高) ，再说方向 (out of the roots) ，最后说结果 (wilts) \par
Watch for:\par
Q2 (a) 常见错误：用 48 除以 20 = 2.4，忘了减去开始的 12 mm\par
Q3 这是 osmosis 换个角度问：水从 high 到 low，所以这次是从根里往外走\par
Script:\par
Q1 她应答：osmosis；xylem。\par
Q2 (a) 让她说算式：48 minus 12 is 36，36 divided by 20 is 1.8 mm per minute。如果她算出 2.4，提醒她先减去开始的 12。\par
Q3 她应答：salty water has a lower concentration of water, so water moves out of the roots by osmosis, and the plant wilts。\par
说：“Q3 是 osmosis 换个角度问，方向还是从高到低，这次是从根里往外。”}
\end{frame}

\begin{frame}{Ways to lose marks: Green Machines}
\relax
\qline[11mm]{\Lref{pollfert}} Pollination moves pollen; fertilisation fuses the nuclei.
\qline[11mm]{\Lref{asexual}} Asexual: one parent, no gametes, identical offspring.
\qline[11mm]{\Lref{adapt}} Feature \textbf{and} how it helps.
\qline[11mm]{\Lref{photoeq}} Light and chlorophyll go above and below the arrow.
\qline[11mm]{\Lref{iodine}} Iodine shows starch, not glucose.
\qline[11mm]{\Lref{xylem}} Xylem carries water; phloem carries glucose.
\fref{Green Machines p.2}
\zh{%
屏幕：Ways to lose marks: Green Machines（植物单元最常见的丢分点）。L1 授粉是送花粉，受精是两个细胞核融合。L2 无性生殖：一个亲本、没有配子、后代一样。L3 写特征也要写它怎么帮助。L4 光和叶绿素写在箭头上下。L6 碘液测淀粉，不是葡萄糖。L7 xylem 运水，phloem 运葡萄糖。(L5 植物一直在呼吸、L8 渗透方向、L9 蒸腾拉力在作业里练。)\par
背景：L 编号是失分点清单，全课统一，Spot the error 的答案也用这些编号。}
\note{%
Teaching line, not a question slide.\par
时间 0:38–0:39。\par
让她每一条说一个今天课上出现过的例子：\Lref{pollfert} Spot the error 1；\Lref{asexual} spider plant；\Lref{photoeq} 和 \Lref{iodine} Spot the error 2；\Lref{xylem} carnation。\par
\Lref{osmosis} 和 \Lref{transp} 在作业 hw1 里还会再练。\par
Script:\par
说：“这六条是今天最容易丢分的地方，你每条说一个刚才出现过的例子。”\par
她可以举：L1 Sam 的第一份答案；L2 spider plant；L4 和 L6 Sam 的第二份答案；L7 康乃馨。}
\end{frame}

\begin{frame}{Summary: Green Machines}
\relax
\vspace{4mm}
{\renewcommand\baselinestretch{1.3}\sSum
Flowers: pollination moves pollen; fertilisation joins the nuclei.\par\vspace{3mm}
Seeds need water, warmth and oxygen; dispersal cuts competition.\par\vspace{3mm}
Leaves: carbon dioxide + water make glucose + oxygen, using light.\par\vspace{3mm}
Water: osmosis into root hairs, up the xylem, out by transpiration.\par}
\fref{Green Machines p.2}
\zh{%
屏幕：Summary: Green Machines（植物单元总结）四句：花：授粉送花粉，受精合细胞核。种子需要水、温暖、氧气；传播减少竞争。叶子：二氧化碳 + 水 用光变成葡萄糖 + 氧气。水：渗透进根毛，沿 xylem 上去，通过蒸腾作用出去。\par
背景：这四句就是整个单元的骨架。}
\note{%
Teaching line, not a question slide.\par
时间 0:39–0:40。\par
让她闭上眼说一遍这四句，再看屏幕对一下。\par
Script:\par
说：“闭上眼睛，把这四句用英文说一遍。”\par
她说完后再看屏幕对一下，漏了哪个词就补哪个词。}
\end{frame}

\begin{frame}{Classwork: Green Machines}
\relax
On your classwork sheet, about 10 minutes:\par\vspace{3mm}
\qline{A} Flowers and seeds
\qline{B} Photosynthesis and respiration
\qline{C} Transport and osmosis, with a graph to draw
\par\vspace{3mm}
Write in full sentences for every \textbf{explain} question.
\fref{Classwork Part A--C}
\zh{%
屏幕：Classwork: Green Machines（课堂练习：植物）。在练习纸上做，大约 10 分钟：A 花和种子；B 光合作用和呼吸作用；C 运输和渗透，有一张图要画。每道 explain 题都写完整句子。\par
背景：答案在 classwork\_answers.pdf。Part C 的图：横轴时间、纵轴高度，0–25 分钟直线上升，之后变平。}
\note{%
Teaching line, not a question slide.\par
时间 0:40–0:50。\par
她做的时候老师不要讲，最后 2 分钟对 Part C 的图 (标尺、轴标签、单位、描点) 。\par
答案在 classwork\_answers.pdf：Part A 重点是 pollination 和 fertilisation 的区分；Part B 重点是 bell jar 实验的 control；Part C 的图从 0 到 25 分钟是直线上升，之后变平。\par
时间不够就把 Part C 的画图留到下课后。\par
Script:\par
说：“现在做 Classwork 的 Part A 到 C，大约 10 分钟，explain 题写完整句子。”\par
她做的时候你不要提示。最后 2 分钟一起检查 Part C 的图：坐标轴、刻度、单位、点有没有描准。}
\end{frame}

% ================================================================ LESSON 2
\begin{frame}
\relax
\lessondivider{LESSON 2 \textperiodcentered{} Metallic materials}{Learning intention: link properties to uses, calculate density, and explain alloys, rusting and the reactions of metals.}
\zh{%
屏幕：LESSON 2 · Metallic materials。学习目标：把金属的性质和用途联系起来，计算 density (密度)，解释 alloys (合金)、rusting (生锈) 和金属的化学反应。\par
背景：这本 booklet 从第 9 页 Alloys 往后她几乎没写，但学校教过。所以这一课每页都先让她说。}
\note{%
Teaching line, not a question slide.\par
时间 0:50–0:51。\par
这本 booklet 从 p.9 Alloys 往后她几乎没写，学校教过。所以 Loop 2–4 每一页都要让她自己说，不要老师先讲。\par
四段：Loop 1 properties and density，Loop 2 heat treatment and alloys，Loop 3 corrosion，Loop 4 reactions and the reactivity series。\par
Script:\par
说：“第二课，metals。这本 booklet 后半本你没写，我们今天一起把它补上。”\par
说：“四段：properties and density；heat treatment and alloys；corrosion (腐蚀)；reactions and the reactivity series (金属活动性顺序)。”\par
Say it:\par
reactivity = ree-ak-TIV-ih-tee\par
alloys = AL-oyz\par
corrosion = kuh-ROH-zhun}
\end{frame}

\begin{frame}{Do Now}
\relax
\qline{1} Give two physical properties shared by most metals.
\qline{2} Name the three common magnetic metals.
\qline{3} Write the formula for density.
\qline{4} Gallium melts at 30\,\textdegree C. Is it solid or liquid on a 35\,\textdegree C day?
\qline{5} A metal is heated red hot, then cooled slowly. Name this heat treatment.
\fref{Metals p.2--8}
\zh{%
屏幕：Do Now。1. 说出大多数金属共有的两个物理性质。2. 说出三种常见的有磁性的金属。3. 写出密度公式。4. 镓 (gallium) 30 °C 熔化，35 °C 的天它是固体还是液体？5. 把金属烧到红热再慢慢冷却，这种热处理叫什么？\par
背景：1 lustrous (有光泽)、malleable (能锤打成形)、ductile (能拉成丝)、sonorous (敲了会响)、good conductor (导热导电好)、熔点高，任选两个。2 iron, nickel, cobalt。3 density = mass ÷ volume。4 液体（温度高于熔点）。5 annealing (退火)。}
\note{%
Q1 any two: lustrous, malleable, ductile, sonorous, good conductors of heat and electricity, high melting point\par
Q2 iron, nickel, cobalt\par
Q3 density $= \dfrac{\text{mass}}{\text{volume}}$\par
Q4 liquid\par
Q5 annealing\par
Short solution:\par
Q4 温度高于 melting point 就是 liquid\par
Q5 annealing = cool slowly；quenching = cool quickly in cold water\par
Watch for:\par
Q1 写 “shiny” 可以，但考试更喜欢 lustrous；写 “strong” 不算通用性质\par
Q3 写成 mass × volume 是 booklet p.4 她表格里出过的错 (\Lref{density})\par
Diagnostic:\par
Q3 错：Loop 1 的 density 两页要慢讲\par
Q5 错：heat treatment 名称不熟，Loop 2 的 Quick review 先做\par
Script:\par
说：“Do Now，五道题，用英文回答。”\par
第 1 题她可以说 lustrous, malleable, ductile, sonorous, good conductor 中任意两个。\par
第 2 题：iron, nickel, cobalt。第 3 题：density = mass divided by volume。如果她说乘，记下来，待会儿重点讲。\par
第 4 题：liquid。第 5 题：annealing。如果她答 quenching (淬火)，说：“quenching 是快速冷却，慢慢冷却叫 annealing。”\par
Say it:\par
malleable = MAL-ee-uh-bul\par
ductile = DUK-tile\par
lustrous = LUS-trus\par
annealing = uh-NEE-ling\par
gallium = GAL-ee-um\par
sonorous = SON-uh-rus\par
quenching = KWEN-ching}
\end{frame}

\begin{frame}{The one idea}
\relax
\vspace{6mm}
\hbox to\textwidth{\hfil\FigMetalIdea\hfil}\par\vspace{7mm}
Top row: physical properties. Bottom row: chemical reactions.\par\vspace{2mm}
Every Metals question sits on one of these two rows.
\fref{Metals p.1}
\zh{%
屏幕：The one idea。上排：原子怎么排列 → 物理性质 → 用途。下排：金属有多活泼 → 它和什么反应 → 怎么保护它。下面两句：上排是物理性质，下排是化学反应；每道金属题都在这两排中的一排。\par
背景：上排包括 malleable、导体、密度、合金；下排包括和氧气、水、酸的反应、生锈和防锈。}
\note{%
Teaching line, not a question slide.\par
时间 0:54–0:55。\par
上排：atoms 怎么排列 → physical properties (malleable, conductor, density) → uses。alloy 也在这一排。\par
下排：reactivity → 和 oxygen、water、acid 的反应 → rusting 和怎么保护。\par
问她：“为什么 gold 做首饰？” 两排都用到：lustrous (上排) + very unreactive，不会 corrode (下排) 。\par
Script:\par
说：“金属这一课分两排：上排是 physical properties，下排是 chemical reactions。”\par
问：“为什么金子用来做首饰？” 她应答：lustrous (上排) + very unreactive, does not corrode (下排) 。}
\end{frame}

% ---------------------------------------------------------- Loop 1
\begin{frame}{Quick review: Property words}
\relax
Match each word to its meaning, then spell it.\par\vspace{2mm}
\begin{tabular}{@{}p{38mm}p{108mm}@{}}
\textbf{1} malleable & \textbf{A} has a shiny surface that reflects light\\
\textbf{2} ductile & \textbf{B} can be hammered into shape\\
\textbf{3} lustrous & \textbf{C} makes a ringing sound when hit\\
\textbf{4} sonorous & \textbf{D} can be stretched into wires\\
\textbf{5} conductor & \textbf{E} lets heat or electric current pass through\\
\end{tabular}
\fref{Metals p.2, p.20}
\zh{%
屏幕：Quick review: Property words（快速回忆：性质词）。配对并拼写：1 malleable 2 ductile 3 lustrous 4 sonorous 5 conductor；A 表面有光泽能反光 B 能锤打成形 C 敲击会发出响声 D 能拉成丝 E 让热或电流通过。\par
背景：答案 1B 2D 3A 4C 5E。拼写难点：malleable 两个 l、-eable 结尾。}
\note{%
1. B\par
2. D\par
3. A\par
4. C\par
5. E\par
Short solution:\par
记法：malleable 像 mallet (锤子) ；ductile 像 duct (管子、拉长) \par
Watch for:\par
1. 拼写：两个 l，-eable 结尾 (malleable) \par
4. sonorous 只有一个 n 在开头\par
Script:\par
说：“配对，然后把 malleable 拼给我听。”\par
她应答：1 B, 2 D, 3 A, 4 C, 5 E。\par
说：“记法：malleable 像 mallet 锤子；ductile 像 duct 管子，拉长。”}
\end{frame}

\begin{frame}{Choosing a metal}
\relax
\renewcommand\arraystretch{1.12}
\begin{tabular}{@{}c c c c c@{}}
\toprule
\textbf{Metal} & \textbf{Melting point} & \textbf{Strength} & \textbf{Electrical} & \textbf{Density}\\
 & \textbf{(\textdegree C)} & & \textbf{conductivity} & \textbf{(g/cm$^{3}$)}\\
\midrule
W & 1450 & high & low & 7.8\\
X & 650 & low & high & 2.7\\
Y & 1080 & medium & very high & 8.9\\
Z & 230 & low & medium & 7.3\\
\bottomrule
\end{tabular}\par\vspace{3mm}
Which metal for (a) a bridge frame, (b) wiring in a house, (c) a light overhead power cable? (d) Which would melt in a Bunsen flame (about 800\,\textdegree C)?
\fref{Metals p.6--7}
\zh{%
屏幕：Choosing a metal（选择金属）。四种虚构金属 W、X、Y、Z 的数据表：熔点、强度、导电性、密度。问：(a) 桥的框架用哪种？(b) 家里的电线用哪种？(c) 要求轻的高架电缆用哪种？(d) 哪几种会在本生灯火焰（约 800 °C）里熔化？\par
背景：(a) W：强度高、熔点高。(b) Y：导电性最好。(c) X：导电好又密度低（2.7），所以轻。(d) X 和 Z：熔点低于 800 °C。}
\note{%
Teaching line, not a question slide.\par
时间 0:57–0:59。\par
(a) W：strength high，melting point high\par
(b) Y：electrical conductivity very high\par
(c) X：conductivity high and density low (2.7)，所以 light\par
(d) X 和 Z：melting point 都低于 800\,\textdegree C\par
这一页先只说选哪个，下一页再练怎么写理由。\par
Script:\par
说：“先只说选哪种，下一页再练怎么写理由。”\par
她应答：(a) W，(b) Y，(c) X，(d) X and Z。\par
如果她 (c) 选 Y，问：“Y 的密度是多少？8.9，比 X 重三倍多。”}
\end{frame}

\begin{frame}{Justify a choice: write it in this order}
\relax
\qline{1} Name the metal.
\qline{2} Give the property that matters for this use.
\qline{3} Say why that property matters for the use.
\qline{4} Compare: why not another metal? Use the data.
\par\vspace{2mm}
{\bfseries A}: 1 + 2 \quad {\bfseries M}: 1 to 3 \quad {\bfseries E}: 1 to 4 with numbers\par\vspace{3mm}
For the power cable I choose \gap[12mm] because it \gap[40mm].\par
This matters because \gap[64mm].\par
Y conducts better, but \gap[64mm].
\fref{Metals p.6--7}
\zh{%
屏幕：Justify a choice: write it in this order（说明选择的理由：按顺序写）。1. 说出金属名。2. 说出这个用途需要的性质。3. 说明这个性质为什么对这个用途重要。4. 比较：为什么不用另一种金属？用数据。A 写 1+2，M 写 1–3，E 写 1–4 并带数字。下面一个填空句：电缆我选\_\_\_\_，因为它\_\_\_\_。这很重要，因为\_\_\_\_。Y 导电更好，但\_\_\_\_。\par
背景：E 级范例：选 X，因为导电性高、密度只有 2.7；轻的电缆给电塔的负担小、不容易下垂；Y 导电更好，但密度 8.9，重三倍多。}
\note{%
Teaching line, not a question slide.\par
时间 0:59–1:01。\par
E 级范例：I choose X because it has a high electrical conductivity and a low density of 2.7 g/cm$^{3}$. A light cable puts less weight on the pylons and does not sag. Y conducts even better, but its density is 8.9 g/cm$^{3}$, more than three times heavier.\par
她 booklet p.6 第 5 题写 “Copper, good conductivity and won't burn you to death”：有 property，但没有 “为什么这个 property 重要”，也没有对比，只有 A (\Lref{justify})。\par
Script:\par
说：“这是选择题写理由的四步，我们用高架电缆做一遍。”\par
她应说出：I choose X because it has high conductivity and a low density of 2.7. A light cable puts less weight on the pylons. Y conducts better, but its density is 8.9, over three times heavier。\par
说：“你 booklet 第 6 页写 copper 时只有性质，没有第 3、4 步，所以只有 A。”}
\end{frame}

\begin{frame}{Density: the formula and the volume}
\relax
\twocol{0.56}{%
density $= \dfrac{\gap[16mm]}{\gap[16mm]}$\par\vspace{2mm}
units: g/cm$^{3}$ or \gap[16mm]\par\vspace{4mm}
Volume of a regular block:\par
\quad \gap[50mm]\par\vspace{2mm}
Volume of an odd-shaped object:\par
\quad the \gap[30mm] method\par\vspace{3mm}
What is the volume of the object on the right?}{0.40}{%
\hbox{\FigCylinder{22}{0}\hspace{8mm}\FigCylinder{29}{1}}\par\vspace{1mm}
\hbox to 40mm{\hspace{2mm}before\hfil after\hspace{6mm}}}
\fref{Metals p.4}
\zh{%
屏幕：Density: the formula and the volume（密度：公式和体积）。左边：density = \_\_\_ ÷ \_\_\_（写成分数）；单位 g/cm$^{3}$ 或 \_\_\_\_；规则物体的体积：\_\_\_\_；不规则物体的体积：\_\_\_\_ 法。右边两个量筒，before 和 after，问：右边物体的体积是多少？\par
背景：density = mass ÷ volume；单位 g/mL（1 cm$^{3}$ = 1 mL，所以两个单位一样）。规则物体：长 × 宽 × 高。不规则物体：displacement (排水法)，体积 = 后读数 − 前读数 = 29 − 22 = 7 mL。}
\note{%
Teaching line, not a question slide.\par
时间 1:01–1:03。\par
空格：mass / volume；g/mL；length × width × height；displacement。\par
问题答案：29 − 22 = 7 mL (读刻度：before 22 mL，after 29 mL) 。\par
1 cm$^{3}$ = 1 mL，所以 g/cm$^{3}$ 和 g/mL 是同一个单位。\par
Script:\par
说：“密度公式，分子分母分别是什么？” 她应答：mass over volume。\par
问单位：g/cm$^{3}$ or g/mL。问规则物体体积：length × width × height。问不规则物体：displacement method。\par
指量筒：“读一下 before 和 after。” 她应答：22 and 29，所以 volume is 7 mL。\par
Say it:\par
displacement = dis-PLAYS-ment}
\end{frame}

\begin{frame}{Density: worked example}
\relax
A brass key has a mass of 25.5\,g. When it is lowered into a measuring cylinder, the water level rises from 30.0\,mL to 33.0\,mL.\par\vspace{3mm}
\qline[14mm]{Step 1} volume $=$ \gap[16mm] $-$ \gap[16mm] $=$ \gap[16mm] mL
\qline[14mm]{Step 2} density $=$ \gap[16mm] $\div$ \gap[16mm]
\qline[14mm]{Step 3} density $=$ \gap[20mm] with its unit
\par\vspace{3mm}
Check: is your answer bigger than 1\,g/mL? Brass sinks, so it should be.
\fref{Metals p.4}
\zh{%
屏幕：Density: worked example（密度例题）。一把黄铜钥匙质量 25.5 g，放进量筒后水面从 30.0 mL 升到 33.0 mL。Step 1 体积 = \_\_\_ − \_\_\_ = \_\_\_ mL；Step 2 密度 = \_\_\_ ÷ \_\_\_；Step 3 密度 = \_\_\_（带单位）。检查：答案比 1 g/mL 大吗？黄铜会沉，所以应该大。\par
背景：体积 33.0 − 30.0 = 3.0 mL；密度 25.5 ÷ 3.0 = 8.5 g/mL。}
\note{%
Teaching line, not a question slide.\par
时间 1:03–1:04。\par
Step 1：$33.0 - 30.0 = 3.0\ \mathrm{mL}$\par
Step 2：$\dfrac{25.5\ \mathrm{g}}{3.0\ \mathrm{mL}}$\par
Step 3：8.5 g/mL\par
\Lref{displace}：volume 是两次读数的差，不是最后的读数。\Lref{density}：除法，单位 g/mL。\par
Script:\par
说：“三步走，你说我听。”\par
她应答：volume 33.0 minus 30.0 equals 3.0 mL；density 25.5 divided by 3.0 equals 8.5 g/mL。\par
说：“最后检查一下：会沉的东西密度一定大于 1，8.5 合理。”}
\end{frame}

\begin{frame}{Spot the error (4)}
\relax
A metal bolt has a mass of 39.5\,g. When it is put into a measuring cylinder, the water rises from 20.0\,mL to 25.0\,mL. Calculate the density of the bolt.\par\vspace{2mm}
\samcard{146mm}{\flawed{volume = 25.0 mL\\density = 39.5 x 25.0 = 987.5 g}}\par\vspace{2mm}
\textbf{Sam's answer has 3 errors.} Find each one, fix it, and say why it loses marks.
\fref{Metals p.4}
\zh{%
屏幕：Spot the error (4)。题目：螺栓质量 39.5 g，放进量筒后水面从 20.0 mL 升到 25.0 mL，求密度。Sam：volume = 25.0 mL；density = 39.5 × 25.0 = 987.5 g。3 个错误。\par
背景：错误 1：体积应是 25.0 − 20.0 = 5.0 mL。错误 2：密度是除法不是乘法。错误 3：单位是 g/mL 不是 g。正确答案 7.9 g/mL（和铁一样）。她 booklet 第 4 页最后一行也犯了乘法的错。}
\note{%
1. volume = 25.0 − 20.0 = 5.0 mL, not 25.0 mL (\Lref{displace})\par
2. density = mass divided by volume, not multiplied (\Lref{density})\par
3. the unit is g/mL, not g (\Lref{density})\par
Short solution:\par
volume $= 25.0 - 20.0 = 5.0\ \mathrm{mL}$；density $= \dfrac{39.5\ \mathrm{g}}{5.0\ \mathrm{mL}} = 7.9$ g/mL (和 booklet p.6 表里的 iron 一样) \par
Watch for:\par
她 booklet p.4 最后一行也是乘了 (16.55 g 和 2 cm 得出 33.25) ，体积单位还写成 cm。答案 987.5 明显不合理，要养成检查 “是否合理” 的习惯。\par
Script:\par
说：“Sam 算密度，3 个错误。”\par
她应找出：volume should be 5.0 mL；divide, not multiply；the unit is g/mL。\par
问：“正确答案是多少？” 39.5 ÷ 5.0 = 7.9 g/mL。\par
说：“987.5 这么大的数一看就不合理，算完一定要看看合不合理。”}
\end{frame}

\begin{frame}{Practice}
\relax
\qline{Q1} A metal block is 2.0\,cm $\times$ 2.0\,cm $\times$ 3.0\,cm and has a mass of 64.8\,g. Calculate its density.\slv{A}
\qline{Q2} A ring has a mass of 31.5\,g. It makes the water in a measuring cylinder rise from 40.0\,mL to 43.0\,mL. Which metal is it?\slv{M}
\par\vspace{1mm}\hspace*{7mm}iron 7.9 \quad copper 8.9 \quad silver 10.5 \quad aluminium 2.7 (g/mL)
\qline{Q3} Silver conducts electricity better than copper. Explain why house wiring is made of copper.\slv{E}
\fref{Metals p.4, p.6--7}
\zh{%
屏幕：Practice。Q1 [A] 金属块 2.0 cm × 2.0 cm × 3.0 cm，质量 64.8 g，求密度。Q2 [M] 戒指质量 31.5 g，量筒水面从 40.0 mL 升到 43.0 mL，是哪种金属？（铁 7.9，铜 8.9，银 10.5，铝 2.7 g/mL）Q3 [E] 银导电比铜好，解释为什么家里的电线用铜。\par
背景：Q1 体积 12.0 cm$^{3}$，密度 5.4 g/cm$^{3}$。Q2 体积 3.0 mL，密度 10.5，是银。Q3 铜导电几乎和银一样好、便宜得多、能拉成丝 (ductile)。}
\note{%
Q1 5.4 g/cm$^{3}$\par
Q2 silver (10.5 g/mL)\par
Q3 copper is almost as good a conductor as silver, is much cheaper, and is ductile, so it can be drawn into wires\par
Short solution:\par
Q1 volume $= 2.0 \times 2.0 \times 3.0 = 12.0\ \mathrm{cm^{3}}$；density $= \dfrac{64.8\ \mathrm{g}}{12.0\ \mathrm{cm^{3}}} = 5.4$ g/cm$^{3}$\par
Q2 volume $= 43.0 - 40.0 = 3.0\ \mathrm{mL}$；density $= \dfrac{31.5\ \mathrm{g}}{3.0\ \mathrm{mL}} = 10.5$ g/mL，对表得 silver\par
Q3 E 要有对比：booklet p.6 的数据 silver 60、copper 50，差别不大，但 silver 是贵金属\par
Watch for:\par
Q2 先减读数再除，答案要和表对上才算数\par
Script:\par
Q1 让她说算式：2.0 × 2.0 × 3.0 = 12.0；64.8 ÷ 12.0 = 5.4 g/cm$^{3}$。\par
Q2：43.0 − 40.0 = 3.0；31.5 ÷ 3.0 = 10.5，是 silver。\par
Q3 她应答：copper is almost as good a conductor, much cheaper, and ductile。要拿 E 一定要和 silver 对比。\par
Say it:\par
aluminium = al-yoo-MIN-ee-um}
\end{frame}

% ---------------------------------------------------------- Loop 2
\begin{frame}{Quick review: Heat treatments}
\relax
\qline{1} Heat red hot, then cool quickly in cold water.
\qline{2} Heat red hot, then cool slowly.
\qline{3} Heat strongly until a blue sheen appears.
\par\vspace{2mm}
Name each treatment: annealing, quenching or tempering.\par
Then match it to a result: \textbf{A} soft, easy to bend \quad \textbf{B} hard but brittle \quad \textbf{C} tough and springy\par\vspace{3mm}
\qline{4} A hair clip must spring back into shape. Which treatment does it need?
\fref{Metals p.8}
\zh{%
屏幕：Quick review: Heat treatments（快速回忆：热处理）。1 烧红后放冷水里快速冷却。2 烧红后慢慢冷却。3 强烈加热直到出现蓝色光泽。说出名称（annealing、quenching、tempering）并配结果：A 软、容易弯；B 硬但脆；C 有韧性、有弹性。4 发夹要能弹回原形，需要哪种处理？\par
背景：1 quenching 淬火 → B；2 annealing 退火 → A；3 tempering 回火 → C；4 tempering。她 booklet 第 8 页把发夹写成 annealing，错了。}
\note{%
1. quenching, B\par
2. annealing, A\par
3. tempering, C\par
4. tempering\par
Short solution:\par
4. hair clip 要 springy，所以是 tempering\par
Watch for:\par
4. 她 booklet p.8 第 2 题写 “Annealing because it modified easily”，这是错的：annealing 让钢变软，夹子会弯了回不去\par
Script:\par
说：“三种热处理，名字加结果。”\par
她应答：1 quenching, B；2 annealing, A；3 tempering, C。\par
问第 4 题：答 tempering，因为要 springy (有弹性)。说：“你 booklet 写的是 annealing，那样夹子会软掉弹不回来。”\par
Say it:\par
tempering = TEM-per-ing}
\end{frame}

\begin{frame}{Why alloys are harder (1)}
\relax
A pure metal: every atom is the same size, in regular layers.\par\vspace{5mm}
\hbox to\textwidth{\hfil\FigPure[5mm]\hspace{14mm}\FigPureSlid[5mm]\hfil}\par\vspace{2mm}
\hbox to\textwidth{\hfil before\hspace{46mm}after a push\hspace{12mm}\hfil}\par\vspace{4mm}
What did the layers do when they were pushed? So is a pure metal hard or soft?
\fref{Metals p.2, p.9}
\zh{%
屏幕：Why alloys are harder (1)（合金为什么更硬 1）。纯金属：所有原子一样大，排成整齐的层。图：左边是原来的层，右边是推一下以后上面两层滑过去了。问：推的时候层怎么样？所以纯金属硬还是软？\par
背景：层很容易互相滑动，所以纯金属软，也因此 malleable、ductile。booklet 说纯铁对很多用途来说太软。}
\note{%
Teaching line, not a question slide.\par
时间 1:07–1:08。\par
答案：the layers slide over each other easily，所以 pure metal 是 soft 的，也因此 malleable 和 ductile。\par
booklet p.9：pure iron is too soft for many uses。\par
Script:\par
说：“这是纯金属的原子，全部一样大，排成一层一层。”\par
问：“What happened when the layers were pushed?” 她应答：the layers slid over each other。\par
问：“So is a pure metal hard or soft?” 答 soft。}
\end{frame}

\begin{frame}{Why alloys are harder (2)}
\relax
An alloy: a metal mixed with another element.\par\vspace{1mm}
\twocol{0.40}{\vspace{1mm}\FigAlloy[4.4mm]}{0.56}{%
\qline{1} The atoms are \gap[28mm] sizes.
\qline{2} This \gap[28mm] the regular layers.
\qline{3} The layers cannot \gap[28mm] easily.
\qline{4} So the alloy is \gap[22mm] than the pure metal.}
\par\vspace{3mm}
{\bfseries A}: harder \quad {\bfseries M}: + different-sized atoms \quad {\bfseries E}: whole chain
\fref{Metals p.9}
\zh{%
屏幕：Why alloys are harder (2)。合金：一种金属和另一种元素混合。图里有大小不同的原子。1. 原子的大小\_\_\_\_。2. 这\_\_\_\_了整齐的层。3. 层不能轻易\_\_\_\_。4. 所以合金比纯金属\_\_\_\_。A：合金更硬；M：加上原子大小不同；E：整条链。\par
背景：different sizes；disrupts；slide over each other；harder。只写“因为是混合物”不给分。}
\note{%
Teaching line, not a question slide.\par
时间 1:08–1:09。\par
空格：different / disrupts / slide over each other / harder。\par
\Lref{alloy}：只写 “because it is a mixture” 不给分，必须写 atoms of different sizes 和 layers can't slide。\par
Script:\par
说：“现在加进一些大原子，看层还能不能滑。”\par
她应答：the atoms are different sizes；this disrupts the layers；the layers cannot slide easily；so the alloy is harder。\par
说：“考试四句都要写，这是 Excellence 的写法。”}
\end{frame}

\begin{frame}{Alloys in use}
\relax
\renewcommand\arraystretch{1.3}
\begin{tabular}{@{}l l l@{}}
\toprule
\textbf{Alloy} & \textbf{Main metals} & \textbf{Use}\\
\midrule
bronze & copper + \gap[18mm] & gears, machine bearings\\
solder & tin + lead & \gap[44mm]\\
stainless steel & iron, chromium, nickel & \gap[44mm]\\
nitinol & nickel + titanium & \gap[44mm]\\
\bottomrule
\end{tabular}\par\vspace{4mm}
Why is solder good for joining electrical wires?\par
Is an alloy a compound or a mixture?
\fref{Metals p.9, p.11}
\zh{%
屏幕：Alloys in use（合金的用途）。表格：bronze (青铜) 铜 + \_\_\_\_，用于齿轮、轴承；solder (焊锡) 锡 + 铅，用于\_\_\_\_；stainless steel (不锈钢) 铁、铬、镍，用于\_\_\_\_；nitinol (镍钛合金) 镍 + 钛，用于\_\_\_\_。问：为什么焊锡适合连接电线？合金是 compound (化合物) 还是 mixture (混合物)？\par
背景：tin（锡）；连接电线；餐具和水槽（不生锈）；牙套和眼镜架（形状记忆，smart alloy 智能合金）。焊锡熔点低，容易熔化，冷了又是固体。合金是混合物：比例可以变，每种成分保持自己的性质；化合物是固定比例、化学结合。}
\note{%
Teaching line, not a question slide.\par
时间 1:09–1:10。\par
空格：tin / joining electrical wires / cutlery and sinks (resists corrosion) / teeth braces and spectacle frames (shape memory)。\par
solder：melting point 低，容易熔化，冷却后在常温下是 solid。\par
alloy 是 mixture：比例可以变，每种成分保留自己的 chemical identity；compound 是固定比例、化学结合。\par
nitinol 是 smart alloy (shape memory alloy)：弯了以后加热或通电会回到原来的形状。\par
Script:\par
说：“填一下表里的空。”\par
她应答：tin；joining electrical wires；cutlery and sinks；teeth braces and glasses frames。\par
问：“Why is solder good for joining wires?” 答 low melting point。问：“Compound or mixture?” 答 mixture。\par
Say it:\par
chromium = KROH-mee-um\par
nitinol = NIT-ih-nol\par
titanium = ty-TAY-nee-um}
\end{frame}

\begin{frame}{Melting points of alloys}
\relax
\twocol{0.46}{\vspace{1mm}\hspace*{6mm}\FigMPGraph[4.9mm]}{0.50}{%
Mixtures of two metals, X and Y.\par\vspace{2mm}
\qline{1} What is the melting point of pure X? Of pure Y?
\qline{2} Which mixture has the lowest melting point?
\qline{3} Describe the trend from 0\,\% to 70\,\% X.}
\fref{Metals p.10--11}
\zh{%
屏幕：Melting points of alloys（合金的熔点）。图：金属 X 和 Y 混合物的熔点，横轴是 X 的百分比，纵轴是熔点。1. 纯 X 的熔点是多少？纯 Y 呢？2. 哪个混合比例熔点最低？3. 描述 0\% 到 70\% X 的趋势。\par
背景：纯 X（100\%）262 °C；纯 Y（0\% X）330 °C；70\% X 时最低，约 190 °C；X 越多熔点越低（到 70\% 为止）。结论：合金的熔点可以比两种纯金属都低。booklet 记图口诀 SLurP of Tea：Scale 刻度、Label 标签、Plot 描点、Title 标题（还有 Axes 坐标轴）。}
\note{%
Teaching line, not a question slide.\par
时间 1:10–1:11。\par
1. pure X (100 \%) 262\,\textdegree C；pure Y (0 \% X) 330\,\textdegree C\par
2. 70 \% X，about 190\,\textdegree C\par
3. the melting point falls steadily as the percentage of X increases\par
提醒 booklet 的 “A SLurP of Tea”：Axes, Scale, Label, Plot, Title。她 p.10 的 solder 图还没画，留在作业里。\par
结论 (和 booklet p.11 Q1 同一个思路) ：alloy 的 melting point 可以比两种纯金属都低。\par
Script:\par
说：“读图题，先看坐标轴。”\par
她应答：pure X 262 °C, pure Y 330 °C；lowest at 70\% X, about 190 °C；the melting point falls as the percentage of X increases。\par
说：“读数时要带单位 °C。”}
\end{frame}

\begin{frame}{Practice}
\relax
\qline{Q1} Name the heat treatment that makes steel hard but brittle.\slv{A}
\qline{Q2} Brass is an alloy of copper and zinc. Explain why brass is harder than pure copper.\slv{M}
\qline{Q3} A car spring and a car door panel are both made of steel. Explain which heat treatment each one should have.\slv{E}
\fref{Metals p.8--9}
\zh{%
屏幕：Practice。Q1 [A] 哪种热处理让钢变硬但变脆？Q2 [M] 黄铜 (brass) 是铜和锌的合金，解释为什么它比纯铜硬。Q3 [E] 汽车的弹簧和车门板都用钢做，各应该做哪种热处理？解释。\par
背景：Q1 quenching。Q2 锌原子和铜原子大小不同，打乱整齐的层，层不能滑动。Q3 弹簧 tempering（有韧性、能弹回）；车门板 annealing（软，容易压成形状）。}
\note{%
Q1 quenching\par
Q2 zinc atoms are a different size from copper atoms; they disrupt the regular layers, so the layers cannot slide over each other easily\par
Q3 spring: tempering, so it is tough and springy and returns to shape; door panel: annealing, so it is soft and easy to press into shape\par
Short solution:\par
Q3 每一个都要 “treatment + 得到的 property + 为什么这个用途需要它”\par
Watch for:\par
Q2 只写 “alloys are harder” 是复述题目，0 分\par
Q3 门板写 quenching 是错的：太 brittle，一撞就裂\par
Script:\par
Q1 她应答：quenching。\par
Q2 她应答：zinc atoms are a different size, they disrupt the layers, so the layers cannot slide easily。\par
Q3 她应答：spring needs tempering so it is springy；door panel needs annealing so it is soft and easy to shape。}
\end{frame}

% ---------------------------------------------------------- Loop 3
\begin{frame}{Quick review: Corrosion and rusting}
\relax
True or false?\par\vspace{2mm}
\qline{1} Rusting is a type of corrosion.
\qline{2} Any metal can rust.
\qline{3} Salt water makes iron rust faster.
\qline{4} Gold corrodes quickly in air.
\qline{5} Corrosion is a chemical reaction of a metal with substances such as oxygen and water.
\fref{Metals p.13, p.21}
\zh{%
屏幕：Quick review: Corrosion and rusting（快速回忆：腐蚀和生锈）。判断对错：1 生锈是腐蚀的一种。2 任何金属都会生锈。3 盐水让铁生锈更快。4 金在空气里很快被腐蚀。5 腐蚀是金属和氧气、水等物质发生的化学反应。\par
背景：1 对；2 错（只有铁和钢叫 rusting，其他金属叫 corrosion）；3 对；4 错（金非常不活泼）；5 对。}
\note{%
1. true\par
2. false\par
3. true\par
4. false\par
5. true\par
Short solution:\par
2. 只有 iron (和 steel) 会 rust；其他金属是 corrode\par
3. booklet p.13：salt solution 和 higher temperature 都会加快 rusting\par
Watch for:\par
LO 要求 “distinguish between corrosion and rusting”：rusting 是 corrosion 的一种，只用于 iron\par
Script:\par
说：“对还是错，用 true 或 false 回答。”\par
她应答：true, false, true, false, true。\par
问第 2 题为什么错：只有 iron 和 steel 会 rust，其他金属是 corrode。}
\end{frame}

\begin{frame}{Rusty nail experiment (1): the set-up}
\relax
\twocol{0.60}{\vspace{2mm}\FigRustTubes[3.1mm]}{0.37}{%
Iron nails, left for 3 days.\par\vspace{2mm}
\qline{A} The calcium chloride removes \gap[24mm].
\qline{B} Boiling removes \gap[24mm] from the water. The oil stops \gap[24mm].
\qline{C} The nail has both \gap[16mm] and \gap[16mm].}
\fref{Metals p.17}
\zh{%
屏幕：Rusty nail experiment (1): the set-up（生锈实验 1：装置）。图：三支试管，每支里一根铁钉，放 3 天。A 底部有 calcium chloride (氯化钙)，塞着塞子；B 是煮过的水，上面一层油，塞着塞子；C 是自来水，开口。填空：A 氯化钙吸走了\_\_\_\_；B 煮沸赶走了水里的\_\_\_\_，油阻止\_\_\_\_；C 里的钉子既有\_\_\_\_又有\_\_\_\_。\par
背景：A 吸走空气里的水分（所以 A 是干燥空气）；B 煮沸赶走溶解在水里的氧气，油阻止空气再溶进水里；C 有水也有氧气。}
\note{%
Teaching line, not a question slide.\par
时间 1:12–1:14。\par
A：water (moisture) from the air，所以 A 是 dry air\par
B：dissolved oxygen (air)；air from dissolving back into the water\par
C：water and oxygen (air)\par
这个设计和 Lesson 1 Practice Q3 的种子实验一样：boiled water + oil = 没有 oxygen。\par
Script:\par
说：“三支试管，我们先搞清楚每支试管缺什么。”\par
她应答：A calcium chloride removes water；B boiling removes oxygen, the oil stops air getting back in；C has both water and oxygen。\par
说：“B 的设计和第一课种子题 D 管一模一样。”\par
Say it:\par
calcium = KAL-see-um}
\end{frame}

\begin{frame}{Rusty nail experiment (2): write it in this order}
\relax
Results after 3 days: \quad A \gap[22mm] \quad B \gap[22mm] \quad C \gap[22mm]\par\vspace{4mm}
\qline{1} Only tube \gap[10mm] rusted.
\qline{2} Tube A had air but no \gap[22mm].
\qline{3} Tube B had water but no \gap[22mm].
\qline{4} So iron needs \gap[50mm] to rust.
\par\vspace{3mm}
{\bfseries A}: C rusted \quad {\bfseries M}: + what A and B lacked \quad {\bfseries E}: the conclusion from all three
\fref{Metals p.17}
\zh{%
屏幕：Rusty nail experiment (2): write it in this order（生锈实验 2：按顺序写）。三天后的结果：A \_\_\_\_ B \_\_\_\_ C \_\_\_\_。1. 只有\_\_\_\_号试管生锈。2. A 管有空气但没有\_\_\_\_。3. B 管有水但没有\_\_\_\_。4. 所以铁生锈需要\_\_\_\_。A：C 生锈；M：再说出 A 和 B 缺什么；E：用三支试管得出结论。\par
背景：A 不生锈，B 不生锈，C 生锈；结论：铁生锈需要 water and oxygen 两样。她 booklet 第 17 页写 B 管“有一点生锈”，这是错的。}
\note{%
Teaching line, not a question slide.\par
时间 1:14–1:15。\par
结果：A no rust；B no rust；C rust。\par
空格：C / water / oxygen / both water and oxygen。\par
她 booklet p.17 的结果写 B “slight rusting, barely can see”，是错的。B 没有 oxygen，不应该生锈；如果真的有一点，说明 boiled water 又溶进了空气 (\Lref{rust})。\par
Script:\par
说：“结果和结论，按四步说。”\par
她应答：A no rust, B no rust, C rust；only C rusted；A had no water；B had no oxygen；so iron needs both water and oxygen to rust。\par
说：“你 booklet 上 B 写了 slight rusting，B 没有氧气，不应该生锈。”}
\end{frame}

\begin{frame}{Spot the error (5)}
\relax
Steel paper clips are sealed in three tubes: P with calcium chloride (dry air), Q in boiled water under a layer of oil, R in tap water with air above. Describe and explain the results.\par\vspace{2mm}
\samcard{146mm}{\flawed{P: no rust, because there is no air. Q: a little rust, because it has water. R: lots of rust. So iron only needs water to rust.}}\par\vspace{2mm}
\textbf{Sam's answer has 3 errors.} Find each one, fix it, and say why it loses marks.
\fref{Metals p.13, p.17}
\zh{%
屏幕：Spot the error (5)。题目：钢制回形针分别封在三支试管里：P 有氯化钙（干燥空气），Q 在煮过的水里、上面有油，R 在自来水里、上面有空气。描述并解释结果。Sam：P 不生锈，因为没有空气。Q 生一点锈，因为有水。R 生很多锈。所以铁生锈只需要水。3 个错误。\par
背景：错误 1：P 有空气，缺的是水。错误 2：Q 不生锈，没有氧气。错误 3：生锈需要水和氧气两样。R 那句是对的。}
\note{%
1. P has air; it does not rust because it has no water (\Lref{rust})\par
2. Q does not rust: boiled water has no oxygen and the oil keeps air out (\Lref{rust})\par
3. iron needs both water and oxygen to rust (\Lref{rust})\par
Short solution:\par
P: no rust, because the calcium chloride removes the water. Q: no rust, because boiling removes the oxygen and the oil keeps air out. R: lots of rust, because it has water and oxygen. So iron needs both water and oxygen to rust.\par
Watch for:\par
R 那一句是对的。steel 也会 rust，因为 steel 主要是 iron。\par
Script:\par
说：“Sam 写的生锈实验，3 个错误。”\par
她应找出：P has air but no water；Q does not rust because there is no oxygen；iron needs both water and oxygen。}
\end{frame}

\begin{frame}{Preventing rust}
\relax
\renewcommand\arraystretch{1.35}
\begin{tabular}{@{}l l l@{}}
\toprule
\textbf{Object} & \textbf{Method} & \textbf{Why this method}\\
\midrule
steel fence & \gap[26mm] & large, cheap, easy to redo\\
door hinge & oil or grease & \gap[40mm]\\
rubbish bin & \gap[26mm] & zinc coating\\
kitchen sink & \gap[26mm] & never needs a coating\\
\bottomrule
\end{tabular}\par\vspace{4mm}
Every method keeps \gap[18mm] and \gap[18mm] away from the iron.
\fref{Metals p.13, p.15}
\zh{%
屏幕：Preventing rust（防止生锈）。表格：钢栅栏 → \_\_\_\_（大、便宜、容易重做）；门铰链 → 上油或润滑脂 → 原因\_\_\_\_；垃圾桶 → \_\_\_\_（镀锌层）；厨房水槽 → \_\_\_\_（永远不需要涂层）。最后一句：每种方法都是让\_\_\_\_和\_\_\_\_接触不到铁。\par
背景：painting 刷漆；铰链会动，漆会被磨掉；galvanising 镀锌；用 stainless steel 不锈钢（合金）。所有方法都是隔开 water 和 oxygen。选哪种方法看物体大小、用途和成本。}
\note{%
Teaching line, not a question slide.\par
时间 1:16–1:17。\par
空格：painting；moving parts，paint would rub off；galvanising；made of stainless steel (alloying)。\par
最后一句：water and oxygen (air)。\par
booklet p.13：选哪种方法看 object 的大小、用途和成本。\par
Script:\par
说：“四种防锈方法，填空。”\par
她应答：painting；moving parts, paint would rub off；galvanising；stainless steel。\par
问最后一句：water and oxygen。\par
Say it:\par
galvanising = GAL-vuh-ny-zing}
\end{frame}

\begin{frame}{Practice}
\relax
\qline{Q1} Name the two substances iron needs in order to rust.\slv{A}
\qline{Q2} Explain why a car in a seaside town rusts faster than the same car far inland.\slv{M}
\qline{Q3} Plan an experiment to find out whether salt water makes iron nails rust faster than tap water. Say what you would change, what you would measure, and what you would keep the same.\slv{E}
\fref{Metals p.13, p.17}
\zh{%
屏幕：Practice。Q1 [A] 铁生锈需要哪两样东西？Q2 [M] 为什么海边小镇的车比内陆的同款车生锈更快？Q3 [E] 设计实验：盐水是否比自来水让铁钉生锈更快？说出你改变什么、测量什么、保持什么不变。\par
背景：Q1 水和氧气。Q2 海边空气里有盐，盐溶液加快生锈。Q3 改变：自来水或盐水；测量：3 天后生锈多少（或钉子质量变化）；不变：同样的钉子（用砂纸打磨干净）、同样体积的水、同样温度、同样时间、都开口；每组做两三次。}
\note{%
Q1 water and oxygen\par
Q2 the air near the sea carries salt, and salt solution speeds up rusting\par
Q3 change: tap water or salt water; measure: how much rust after 3 days; keep the same: type and size of nail (cleaned with sandpaper), volume of water, temperature, time, tubes open to the air; repeat each tube\par
Short solution:\par
Q3 E 要有：independent variable、dependent variable、至少三个 controlled variables、重复实验\par
Watch for:\par
Q3 她实验题最容易漏 “how to measure”：写 “compare the amount of rust” 或 “measure the mass of the nail before and after” 都可以\par
Script:\par
Q1 她应答：water and oxygen。\par
Q2 她应答：the air by the sea carries salt, and salt speeds up rusting。\par
Q3 让她说出三样：change tap water or salt water；measure the amount of rust after 3 days；keep the same nail, volume of water, temperature, time；repeat。\par
说：“实验设计题最容易漏的是 ‘怎么测量’。”}
\end{frame}

% ---------------------------------------------------------- Loop 4
\begin{frame}{Quick review: Word equation patterns}
\relax
\qline{1} metal + oxygen $\to$ \gap[40mm]
\qline{2} metal + water $\to$ \gap[34mm] + \gap[26mm]
\qline{3} metal + hydrochloric acid $\to$ \gap[34mm] + \gap[26mm]
\qline{4} Test for hydrogen: a lit splint gives a \gap[30mm].
\fref{Metals p.12, p.18--19}
\zh{%
屏幕：Quick review: Word equation patterns（快速回忆：文字方程式规律）。1 金属 + 氧气 → \_\_\_\_。2 金属 + 水 → \_\_\_\_ + \_\_\_\_。3 金属 + 盐酸 (hydrochloric acid) → \_\_\_\_ + \_\_\_\_。4 检验氢气：点燃的木条发出\_\_\_\_。\par
背景：1 metal oxide 金属氧化物；2 metal hydroxide 金属氢氧化物 + hydrogen 氢气；3 metal chloride 金属氯化物 + hydrogen；4 squeaky pop（“噗”的一声）。}
\note{%
1. metal oxide\par
2. metal hydroxide + hydrogen\par
3. metal chloride + hydrogen\par
4. squeaky pop\par
Short solution:\par
2 和 3 都放出 hydrogen gas，所以都可以用 lit splint 测\par
Watch for:\par
2. 写成 metal oxide + hydrogen 是常见错误 (\Lref{wordeq})\par
Script:\par
说：“三个规律记住了，所有方程式都能写。”\par
她应答：metal oxide；metal hydroxide + hydrogen；metal chloride + hydrogen；squeaky pop。\par
说：“金属和水生成的是 hydroxide，不是 oxide，这是常见错误。”\par
Say it:\par
hydroxide = hy-DROK-side\par
hydrochloric = hy-droh-KLOR-ik}
\end{frame}

\begin{frame}{Writing word equations}
\relax
\qline{(a)} Iron wool burns in oxygen.\par\hspace*{7mm}iron + oxygen $\to$ \gap[40mm]
\qline{(b)} Lithium fizzes in cold water.\par\hspace*{7mm}lithium + water $\to$ \gap[40mm] + \gap[22mm]
\qline{(c)} Aluminium powder in dilute hydrochloric acid.\par\hspace*{7mm}\gap[120mm]
\qline{(d)} Gold in dilute hydrochloric acid.\par\hspace*{7mm}\gap[60mm]
\fref{Metals p.12, p.18--19}
\zh{%
屏幕：Writing word equations（写文字方程式）。(a) 钢丝绒在氧气里燃烧：iron + oxygen → \_\_\_\_。(b) 锂在冷水里冒泡：lithium + water → \_\_\_\_ + \_\_\_\_。(c) 铝粉放进稀盐酸：写出整个方程式。(d) 金放进稀盐酸：\_\_\_\_。\par
背景：(a) iron oxide；(b) lithium hydroxide + hydrogen；(c) aluminium + hydrochloric acid → aluminium chloride + hydrogen；(d) no reaction（金在活动性顺序最底下）。}
\note{%
Teaching line, not a question slide.\par
时间 1:21–1:23。\par
(a) iron oxide\par
(b) lithium hydroxide + hydrogen\par
(c) aluminium + hydrochloric acid → aluminium chloride + hydrogen\par
(d) no reaction：gold 在 reactivity series 的最下面\par
考试写 word equation：用词，不用化学式；箭头不要写成等号。\par
Script:\par
说：“用刚才的规律写四个方程式，用词，不用化学式。”\par
她应答：(a) iron oxide；(b) lithium hydroxide plus hydrogen；(c) aluminium chloride plus hydrogen；(d) no reaction。\par
Say it:\par
lithium = LITH-ee-um}
\end{frame}

\begin{frame}{The reactivity series}
\relax
\vspace{3mm}
\hbox to\textwidth{\hfil\FigReactivity\hfil}\par\vspace{4mm}
\qline{1} Would tin react with dilute acid?
\qline{2} Why is gold found in the ground as pure metal, but iron is not?
\fref{Metals p.20}
\zh{%
屏幕：The reactivity series（金属活动性顺序）。图：从最活泼到最不活泼分四组：钾、钠、锂、钙 | 镁、铝、锌、铁、锡、铅 | 铜、汞、银 | 金。下面三条箭头：和水反应（第一组）、和酸反应（第一、二组）、和氧气反应（第一到三组）。问：1. 锡会和稀酸反应吗？2. 为什么地里能找到纯金，却找不到纯铁？\par
背景：1 会（锡在“和酸反应”的范围里，只是很慢）。2 金非常不活泼，不和氧气、硫反应，所以以纯金属存在；铁比较活泼，以化合物的形式存在于 ores (矿石) 里，要 extracted (提炼)。}
\note{%
Teaching line, not a question slide.\par
时间 1:23–1:25。\par
1. yes：tin 在 “react with acids” 的范围里 (反应很慢) \par
2. gold is so unreactive that it does not react with oxygen or sulfur, so it stays as the pure metal; iron is more reactive, so it is found as compounds (ores) and must be extracted\par
这一页对应 booklet p.20 Summary 的图，重画成四组。\par
Script:\par
说：“这是 booklet 第 20 页的图，越往左越活泼。”\par
问第 1 题：答 yes，tin is in the ‘react with acids’ range。\par
问第 2 题：她应答：gold is very unreactive so it stays as the metal；iron is more reactive so it is found as compounds in ores。\par
Say it:\par
potassium = puh-TAS-ee-um\par
magnesium = mag-NEE-zee-um\par
mercury = MER-kyuh-ree}
\end{frame}

\begin{frame}{Spot the error (6)}
\relax
Complete each word equation, or write ``no reaction''.\par\vspace{2mm}
\samcard{146mm}{\flawed{1. calcium + oxygen makes calcium oxygen\\
2. iron + hydrochloric acid makes iron chloride + water\\
3. potassium + water makes potassium hydroxide + hydrogen\\
4. silver + hydrochloric acid makes silver chloride + hydrogen}}\par\vspace{2mm}
\textbf{Sam's answer has 3 errors.} Find each one, fix it, and say why it loses marks.
\fref{Metals p.12, p.18--20}
\zh{%
屏幕：Spot the error (6)。题目：完成文字方程式，或写“no reaction”。Sam：1 钙 + 氧气 → calcium oxygen；2 铁 + 盐酸 → 氯化铁 + 水；3 钾 + 水 → 氢氧化钾 + 氢气；4 银 + 盐酸 → 氯化银 + 氢气。3 个错误。\par
背景：错误 1：应是 calcium oxide。错误 2：金属和酸生成 hydrogen 不是 water。错误 3（第 4 条）：银不和稀酸反应，no reaction。第 3 条是对的。}
\note{%
1. calcium + oxygen → calcium oxide, not calcium oxygen (\Lref{wordeq})\par
2. iron + hydrochloric acid → iron chloride + hydrogen, not water (\Lref{wordeq})\par
3. silver + hydrochloric acid → no reaction: silver is below the acids in the series (\Lref{react})\par
Short solution:\par
calcium + oxygen → calcium oxide; iron + hydrochloric acid → iron chloride + hydrogen; potassium + water → potassium hydroxide + hydrogen (correct); silver + hydrochloric acid → no reaction\par
Watch for:\par
第 3 条是对的，error 的编号按找到的顺序 1–3，不按 Sam 的行号。\par
Script:\par
说：“四条方程式，3 条有错，找出来。”\par
她应找出：calcium oxide, not calcium oxygen；iron chloride plus hydrogen, not water；silver gives no reaction。\par
问：“第 3 条对吗？” 答：对。}
\end{frame}

\begin{frame}{Practice}
\relax
\qline{Q1} Write the word equation for aluminium reacting with oxygen.\slv{A}
\qline{Q2} Four metals are added to dilute acid. Put them in order of reactivity, most reactive first, and give your evidence.\slv{M}
\par\vspace{1mm}\hspace*{7mm}\begin{tabular}{@{}l l@{}}
P & many bubbles, quickly\\
Q & no bubbles\\
R & a few bubbles, slowly\\
S & reacts even with cold water\\
\end{tabular}
\qline{Q3} Describe an experiment with dilute hydrochloric acid to show that zinc is more reactive than copper. Give the result you expect.\slv{E}
\fref{Metals p.18--20}
\zh{%
屏幕：Practice。Q1 [A] 写出铝和氧气反应的文字方程式。Q2 [M] 四种金属放进稀酸：P 气泡多又快；Q 没有气泡；R 气泡少又慢；S 连冷水都能反应。按活泼程度从高到低排列并给证据。Q3 [E] 描述一个用稀盐酸证明锌比铜活泼的实验，并说出预期结果。\par
背景：Q1 aluminium + oxygen → aluminium oxide。Q2 S, P, R, Q。Q3 同样大小的锌片和铜片放进同样体积、同样浓度的稀盐酸：锌冒泡（氢气），铜不反应。}
\note{%
Q1 aluminium + oxygen → aluminium oxide\par
Q2 S, P, R, Q\par
Q3 add same-sized pieces of zinc and copper to the same volume of the same dilute hydrochloric acid; zinc fizzes (hydrogen), copper does not react\par
Short solution:\par
Q2 evidence：S 连冷水都能反应，最活泼；P 气泡多又快；R 少又慢；Q 没有反应\par
Q3 fair test：same size, same volume, same concentration, same temperature；可以用 lit splint 证明气体是 hydrogen\par
Watch for:\par
Q2 只排顺序不写证据只有 A (\Lref{react})\par
Script:\par
Q1 她应答：aluminium plus oxygen gives aluminium oxide。\par
Q2 她应答：S, P, R, Q，并说证据：S reacts even with cold water, P bubbles fastest, R slowly, Q not at all。\par
Q3 她应答：same size pieces, same volume and concentration of acid；zinc fizzes, copper does not react。}
\end{frame}

\begin{frame}{Ways to lose marks: Metals}
\relax
\qline[11mm]{\Lref{density}} Density = mass $\div$ volume, unit g/mL or g/cm$^{3}$.
\qline[11mm]{\Lref{displace}} Volume = final reading $-$ first reading.
\qline[11mm]{\Lref{justify}} Property + why it matters + comparison.
\qline[11mm]{\Lref{alloy}} Different-sized atoms, layers cannot slide.
\qline[11mm]{\Lref{rust}} Rusting needs water \textbf{and} oxygen.
\qline[11mm]{\Lref{wordeq}} Oxide; hydroxide + hydrogen; chloride + hydrogen.
\fref{Metals p.1}
\zh{%
屏幕：Ways to lose marks: Metals（金属单元最常见的丢分点）。L10 密度 = 质量 ÷ 体积，单位 g/mL 或 g/cm$^{3}$。L11 体积 = 后读数 − 前读数。L12 性质 + 为什么重要 + 对比。L13 原子大小不同，层不能滑动。L14 生锈需要水和氧气。L15 氧化物；氢氧化物 + 氢气；氯化物 + 氢气。\par
背景：这些编号接着植物单元的 L1–L9 往下排。L16（活动性要有证据、铜银金不和酸反应）在作业里练。}
\note{%
Teaching line, not a question slide.\par
时间 1:27–1:28。\par
\Lref{density} 和 \Lref{displace} 来自她 booklet p.4 的表；\Lref{rust} 来自 p.17 的 tube B；\Lref{justify} 来自 p.6。\par
\Lref{react} 在作业 hw2 里再练。\par
Script:\par
说：“金属这六条，每条说一个刚才的例子。”\par
她可以举：L10、L11 Sam 的螺栓；L13 黄铜；L14 回形针；L15 Sam 的方程式。}
\end{frame}

\begin{frame}{Summary: Metallic materials}
\relax
\vspace{4mm}
{\renewcommand\baselinestretch{1.3}\sSum
Properties decide uses: say the property, why it matters, and compare.\par\vspace{3mm}
Density = mass $\div$ volume; displacement gives the volume.\par\vspace{3mm}
Alloys are harder: different-sized atoms stop the layers sliding.\par\vspace{3mm}
Rust needs water and oxygen; reactivity decides how metals react.\par}
\fref{Metals p.20}
\zh{%
屏幕：Summary: Metallic materials（金属单元总结）四句：性质决定用途：说出性质、为什么重要、并比较。密度 = 质量 ÷ 体积；排水法测体积。合金更硬：不同大小的原子阻止层滑动。生锈需要水和氧气；活泼程度决定金属怎么反应。\par
背景：同样，这四句是整个单元的骨架。}
\note{%
Teaching line, not a question slide.\par
时间 1:28–1:29。\par
同样让她先说再看。\par
Script:\par
说：“还是先闭眼说一遍，再看屏幕对一下。”}
\end{frame}

\begin{frame}{Classwork: Metallic materials}
\relax
On your classwork sheet, about 10 minutes:\par\vspace{3mm}
\qline{D} Properties and density
\qline{E} Alloys and heat treatment
\qline{F} Corrosion and reactions, with an experiment to plan
\par\vspace{3mm}
Show every step of a calculation, with units.
\fref{Classwork Part D--F}
\zh{%
屏幕：Classwork: Metallic materials（课堂练习：金属）。大约 10 分钟：D 性质和密度；E 合金和热处理；F 腐蚀和反应，有一个实验要设计。计算题每一步都写出来，带单位。\par
背景：答案在 classwork\_answers.pdf。Part D 密度题答案 2.7 g/cm$^{3}$（铝）；Part F 实验设计要有三个不变量。}
\note{%
Teaching line, not a question slide.\par
时间 1:29–1:40。\par
答案在 classwork\_answers.pdf：Part D 的 density 题答案是 2.7 g/cm$^{3}$ (aluminium)；Part E 重点是 alloy 的解释链；Part F 的实验设计要有三个 controlled variables。\par
1:38 收卷，准备发 Mini mock。\par
Script:\par
说：“做 Classwork 的 Part D 到 F，计算要写步骤和单位。”\par
1:38 收卷，准备发 mini mock。}
\end{frame}

% ================================================================ close
\begin{frame}{Mini mock}
\relax
\vspace{4mm}
{\sEmph\bfseries 17 minutes \textperiodcentered{} both units \textperiodcentered{} on the paper}\par\vspace{6mm}
Answer every part.\par\vspace{2mm}
Each question is marked A, M or E.\par\vspace{2mm}
Explain questions: write it in the order you practised today.
\fref{Mini mock}
\zh{%
屏幕：Mini mock（小模考）。17 分钟 · 两个单元 · 写在卷子上。每一问都要答。每题按 A、M、E 评分。Explain 题按今天练习的顺序写。\par
背景：发 mock.pdf，计时 17 分钟，不提示。答案和评分标准在 mock\_answers.pdf。}
\note{%
Teaching line, not a question slide.\par
时间 1:40–1:57。\par
发 mock.pdf，计时 17 分钟，老师不提示。\par
答案和评分标准在 mock\_answers.pdf，每题写了 A、M、E 各需要写到哪一步。\par
时间到了先收卷，课后批改，下次课开头花 5 分钟讲最差的两题。\par
Script:\par
说：“现在是 17 分钟的小模考，两个单元都有，我不提示。explain 题用今天练的顺序写。”\par
时间到了先收卷，课后批改，下次课开头花 5 分钟讲最差的两题。}
\end{frame}

\begin{frame}{Homework}
\relax
\qline{hw1} Green Machines, about 40 minutes
\qline{hw2} Metallic materials, about 40 minutes
\par\vspace{3mm}
Also: Metals booklet p.10--11, draw the solder graph and answer Q1--2.\par\vspace{3mm}
Do them on two different days before the exam.
\fref{hw1, hw2}
\zh{%
屏幕：Homework（作业）。hw1 Green Machines，约 40 分钟；hw2 Metallic materials，约 40 分钟。另外：Metals booklet 第 10–11 页，画焊锡的图并回答 Q1–2。分两天、在考试前做完。\par
背景：答案文件 hw1\_answers.pdf 和 hw2\_answers.pdf，同版面红字。}
\note{%
Teaching line, not a question slide.\par
时间 1:57–2:00。\par
hw1：pollination、osmosis、transpiration 的 Explain 题和一个 Spot the error；hw2：density 计算、reactivity 数据题、rusting 实验设计和一个 Spot the error。\par
答案文件 hw1\_answers.pdf 和 hw2\_answers.pdf，同版面红色答案。\par
booklet p.10 的 solder 图：提醒 SLurP of Tea，x 轴是 percentage of tin。\par
Script:\par
说：“两份作业，各 40 分钟，分两天做。另外把 Metals booklet 第 10 页的焊锡图画了。”\par
提醒她画图用 SLurP of Tea：坐标轴、刻度、标签、描点、标题，x 轴是 percentage of tin。}
\end{frame}

\end{document}
